\documentclass[11pt]{article}
\usepackage[margin=1in]{geometry}
\usepackage[T1]{fontenc}
\usepackage{lmodern,microtype,amsmath,amssymb,bm,booktabs,graphicx}
\usepackage[round,authoryear]{natbib}
\usepackage{xurl}
\usepackage[colorlinks=true,allcolors=blue]{hyperref}
\usepackage{caption,subcaption}
\usepackage{float}
\floatstyle{plaintop}
\restylefloat{table}
\usepackage{longtable}
\usepackage{placeins}
\usepackage{siunitx}
\graphicspath{{figures/}}
\newcommand{\Ua}{U_{\mathrm a}}
\newcommand{\Ub}{U_{\mathrm b}}
\newcommand{\Cp}{C_{\mathrm p}}
\newcommand{\CpH}{C_{\mathrm{p,H}}}
\newcommand{\CpS}{C_{\mathrm{p,S}}}
\newcommand{\Tamb}{T_{\mathrm{amb}}}
\newcommand{\draftstatus}{Draft 0 --- AI drafted; awaiting human scientific review}
\title{Distributed-delay formulations of a thermal laboratory model:\\
identifiability, uncertainty, and experiment design}
\author{Research working draft}
\date{October 21, 2026}
\begin{document}
\maketitle
\begin{center}\small\textbf{\draftstatus}\end{center}
\begin{abstract}

% Expanded input: sections/abstract
The Temperature Control Laboratory (TC-Lab) provides transient heater and sensor
data for studying parameter estimation and experiment design. We compared
$n$-stage input-delay models in conductance--capacity and gain--time-constant
coordinates, separating formulation, initialization and numerical propagation.
The original four-parameter thermal model has an exact one-dimensional physical
parameter equivalence family under the modeled output and prescribed initial
conditions. The cold one-stage formulation represents its transfer behavior
with three parameters and a global pole-label ambiguity. Constant input with
equilibrated delay states leaves delay unobservable. Coordinate changes alter
conditioning without adding information. Cold joint calibration favors one
stage, while initialization changes the order ranking and transfer errors on
held-out records exceed calibration errors. Profiles, resampling and controlled
coverage comparisons show strong sensitivity to serial-error assumptions;
HAC intervals improve but do not restore nominal coverage in the longer
correlation scenario. Ambient/history and command-hold sensitivity further
limit the interpretation of narrow conditional intervals. Constrained local
input designs improve their declared information and fixed-mean separation
criteria. Refitting competitors removes equilibrated constant-input separation
and qualifies fixed-mean discrimination claims. The results support
identifiable transfer functions and explicit experimental assumptions as the
basis for inference. This draft and all scientific choices await human review.

% End expanded input

\end{abstract}

% Expanded input: sections/introduction
\section{Introduction and theoretical grounding}
Parameter estimation for the TC-Lab must distinguish uncertainty in a measured
thermal response from ambiguity in its physical explanation. A model may predict
the observed temperature closely while several parameter combinations produce
the same response. Changing parameter coordinates can improve numerical scaling,
but an invertible coordinate transformation preserves that ambiguity.

Continuous-time identification also distinguishes the input-driven model from
the residual process. The latter can include measurement error, unmeasured
disturbances, and model inadequacy \citep{garnier-2008-direct}. A fitted residual
covariance therefore does not by itself identify sensor noise or establish a
physical disturbance mechanism. This distinction guides the uncertainty
comparisons below: we vary working noise assumptions and examine the errors
that remain on experiments excluded from calibration.

Identifiability concerns the uniqueness of parameters given a model, an observed
output, and an experiment. Practical estimability additionally depends on the
record length, input, noise, and initial-state information. Profile likelihood
examines a parameter while reoptimizing the remaining parameters and can reveal
features missed by a local quadratic approximation
\citep{raue-2009-structural}. We combine that approach with explicit transfer
functions, sensitivity singular values, and resampling. The algebraic calculation
is needed to distinguish an exact equivalence class from a numerically shallow
objective or a finite exploration bound.

Linear time-invariant system identification supplies a connection between input
excitation, parameter information, and prediction accuracy. Application-oriented
design can target informative parameter combinations rather than require equal
precision in every physical parameter. \citet{hjalmarsson-2009-system} develops
this distinction using parameter confidence regions and application loss. Its
prediction-error information formula assumes that the true system belongs to
the model set and that the predictor is differentiable. Those assumptions require
separate scrutiny here because the TC-Lab residuals can contain serial structure
and between-experiment discrepancies.

Finite-horizon amplitude constraints introduce a further distinction between an
information objective and its numerical solution. \citet{manchester-2010-input}
formulates time-domain information design as a nonconvex input problem and
develops a convex relaxation. We use bounded multistart nonlinear optimization
with explicit command and temperature constraints. We have not implemented that
relaxation, so its approximation guarantees do not apply to the designs reported
here. Comparing local solutions across assumed models and residual covariances
is part of assessing a proposed experiment.

An $n$-stage chain is a family of distributed-delay models. The linear chain
construction relates cascaded first-order stages to Erlang dwell times
\citep{hurtado-2019-generalizations}. For a fixed mean delay, increasing the stage
count narrows the delay distribution and approaches a pure delay. Fitting this
family therefore tests a shape assumption as well as a numerical approximation.
Its auxiliary stage count differs from its number of estimated parameters.

This study asks how these distinctions affect the available TC-Lab data. We
compare physical and input--output parameterizations, quantify conditional
uncertainty, test sensitivity to initialization and data handling, and evaluate
simulated inputs for future estimation and discrimination. The study is a
reproducible draft for internal review; it does not establish a validated
physical mechanism from temperature fit alone.

Appendix~\ref{sec:frequency-information} derives the frequency sensitivities
of the delayed models and an ideal periodic information matrix. It connects the
low-frequency lag ambiguity and repeated-pole singularity to the information
criteria used in the finite-horizon designs.

% End expanded input


% Expanded input: sections/models
\section{Model formulations}
\label{sec:models}

\subsection{Thermal balance with a distributed input delay}
The recorded input $u(t)$ is the heater command in percentage points, and the measured
output $T(t)$ is the first temperature sensor in degrees Celsius. We fixed the
heater conversion $c=\alpha P$ at $0.032\,\mathrm{W}$ per percentage point, using
the inherited values $\alpha=0.00016$ and $P=200$. The ambient temperature
$\Tamb$ was set to the first measured temperature in each experiment. This
choice conditions the analysis on an uncertain ambient estimate.

For a chain of $n$ delay states, the auxiliary variables $z_i(t)$ have the same
units as $u(t)$. The physical parameters $\Ua>0$, $\Cp>0$, and $\theta>0$
denote thermal conductance, effective heat capacity, and mean input-delay time,
respectively. The model is
\begin{align}
 \dot z_1 &= \frac{n}{\theta}(u-z_1), &
 \dot z_i &= \frac{n}{\theta}(z_{i-1}-z_i),\quad i=2,\ldots,n,\label{eq:chain}\\
 \Cp\dot T &= -\Ua(T-\Tamb)+c z_n.\label{eq:physical}
\end{align}
The input--output (IO) coordinates $K=c/\Ua$ and $\tau=\Cp/\Ua$ are steady-state
gain and thermal time constant. In these coordinates the temperature equation is
\begin{equation}
 \dot T=-\frac{T-\Tamb}{\tau}+\frac{K}{\tau}z_n.
 \label{eq:io}
\end{equation}
Both parameterizations describe the same positive model family. Their inverse
map is $\Ua=c/K$ and $\Cp=c\tau/K$. Distinct rectangular parameter bounds need
not map onto each other; our numerical comparison records active bounds and
compares common interior solutions.

For zero initial deviations, the Laplace variable $s$ gives the transfer function
\begin{equation}
 G_n(s)=\frac{K}{(1+\tau s)(1+\theta s/n)^n}.
 \label{eq:transfer}
\end{equation}
The delay chain has an Erlang impulse-response kernel with mean $\theta$ and
variance $\theta^2/n$; Appendix~\ref{sec:structural} derives these quantities.
The linear chain construction is consistent with the Erlang dwell-time
representation studied by \citet{hurtado-2019-generalizations}. A finite chain
responds before time $\theta$ and therefore represents a distributed lag. As
$n$ increases at fixed $\theta$, the delay factor $(1+\theta s/n)^{-n}$ tends to
$\exp(-\theta s)$. We also fitted the limiting pure-delay model,
\begin{equation}
 \tau\dot T(t)=-(T(t)-\Tamb)+K u(t-\theta).
 \label{eq:puredelay}
\end{equation}
Increasing $n$ changes the transfer-function shape and adds dynamic states;
it does not add estimated parameters to the three-parameter family.

\subsection{Initial-state assumptions}
We conditioned $T(0)$ on the first observation. The cold-start assumption set
all $z_i(0)=0$, while the input-equilibrated assumption set all $z_i(0)=u(0)$.
The pure-delay model used the corresponding zero or constant input prehistory.
These assumptions encode different unobserved histories of heater operation.
Neither was established by the recorded input before the first sample, because
that history is absent from the supplied data.
Input-equilibrated refers to delay memory. Under a constant command $u_0$,
equilibrium of the entire model would additionally require
$T(0)=\Tamb+Ku_0$. Here the thermal initial condition is fixed independently
at the first reading. Determining the preceding thermal trajectory requires
pre-window observations.

\subsection{Original two-state thermal model}
The inherited model describes heater temperature $T_{\mathrm H}$ and sensor
temperature $T_{\mathrm S}$. Its conductances $\Ua,\Ub>0$ and heat capacities
$\CpH,\CpS>0$ satisfy
\begin{align}
 \CpH\dot T_{\mathrm H}&=-\Ua(T_{\mathrm H}-\Tamb)
                       -\Ub(T_{\mathrm H}-T_{\mathrm S})+cu,\\
 \CpS\dot T_{\mathrm S}&=\Ub(T_{\mathrm H}-T_{\mathrm S}).
 \label{eq:original}
\end{align}
Only $T_{\mathrm S}$ was used in this single-output comparison. Both initial
temperatures were fixed at the first sensor reading. Under these conditions its
transfer function depends on three combinations of four physical parameters:
\begin{equation}
 G_{\mathrm{orig}}(s)=\frac{K}{1+a s+b s^2},\qquad
 K=\frac{c}{\Ua},\quad
 a=\frac{\CpH}{\Ua}+\frac{\CpS}{\Ub}+\frac{\CpS}{\Ua},\quad
 b=\frac{\CpH\CpS}{\Ua\Ub}.
 \label{eq:original-transfer}
\end{equation}
Appendix~\ref{sec:structural} exhibits the resulting continuum of equivalent
physical parameters. The one-stage cold-start model has $a=\tau+\theta$ and
$b=\tau\theta$, so its zero-state transfer family coincides with this original
model when the poles are distinct. The one-stage model also permits a repeated
pole, which is a limiting case of the strictly positive original family:
the original family has $a^2>4b$, whereas the repeated pole has $a^2=4b$.
The pole labels $\tau$ and $\theta$ can be interchanged. We used the
convention $\tau\geq\theta$ when the exchanged point remained within the
specified parameter bounds and retained both modes in likelihood profiles.

% End expanded input


% Expanded input: sections/methods
\section{Computational methods}
\label{sec:methods}

\subsection{Data and reproducibility}
We preserved the supplied observations and first reproduced the unchanged
inherited simulations twice. The resulting trajectories were byte-identical
within each dataset. We then exported the teaching notebooks to Python scripts
with cell markers and source hashes, checking equality of their scientific
abstract syntax trees. The delay study uses separate Python modules for data,
simulation, inference, and result generation.

The data audit found 15 CSV files. Two validation files were numeric duplicates
of the sine and step calibration files, respectively. We excluded both duplicates
from validation, leaving 11 files. We fitted the sine experiment, the step
experiment, and their combination separately. Each calibration experiment has
901 temperature observations. The first observation conditions the initial
temperature and is excluded from the fitted residual vector, leaving 900
residuals per experiment. This treatment does not propagate uncertainty in the
initial temperature or in the ambient temperature inferred from it.

Long calculations run from immutable copies of committed source and raw data.
The run manifests retain source, configuration and data hashes, package versions,
random seeds, every optimizer attempt, and failures. A single manifest identifies
each numerical phase. Result tables and figures are generated from saved outputs.

\subsection{Simulation and fitting}
We used exact state propagation of the linear time-invariant dynamics over each
sampling interval, with linear interpolation of the recorded heater command.
Separate checks compared left and right input holds and the inherited backward
Euler discretization. The pure-delay simulator integrates across shifted input
knots. These choices distinguish input reconstruction and numerical discretization
from a change in model equations.

The baseline comparison fitted seven orders, $n\in\{1,2,3,4,6,8,12\}$, two
initial-state assumptions, and both parameter coordinates. Pure-delay fits used
the same two initial-state assumptions and coordinates. The original model used
its four physical parameters. These choices gave 99 calibration cases.
We minimized unweighted temperature sum of squared errors (SSE) using bounded
least squares in logarithmic parameters, with 24 starts per case and seed 532.
The broad input--output bounds were $0.2\leq K\leq1.5$,
$5\leq\tau\leq600\,\mathrm{s}$, and
$0.01\leq\theta\leq200\,\mathrm{s}$. The physical-coordinate bounds were
$0.032/1.5\leq\Ua\leq0.032/0.2\,\mathrm{W\,K^{-1}}$,
$0.1\leq\Cp\leq100\,\mathrm{J\,K^{-1}}$, and the same delay bounds.
The original model used $\Ua\in[0.005,0.3]$ and
$\Ub\in[0.0001,1]$ in $\mathrm{W\,K^{-1}}$, with
$\CpH\in[0.01,100]$ and $\CpS\in[0.001,20]$ in
$\mathrm{J\,K^{-1}}$. The reduced physical rectangle is not the exact image
of the IO rectangle. Same-point covariance and information comparisons separate
coordinate effects from changes in the feasible parameter domain.
We retained convergence status and active bounds rather than interpreting every
optimizer return as a successful estimate.
Root mean squared error (RMSE) uses the same retained residuals.

\subsection{Ambient, history and input-reconstruction sensitivity}
We prespecified 241 additional calibrations. For each of eight delayed models
and three calibration datasets, four nuisance treatments give 96 fits:
free ambient with cold history, free ambient with equilibrated history,
free history with fixed ambient, and both free. Ambient offsets have bounds
$[-2,2]\,{}^\circ\mathrm C$, separately for each record; history fractions $\gamma_j$
have bounds $[0,1]$. These bounds describe assumed sensitivity ranges.
They are not independently measured probability distributions. The three
IO parameters are shared across records in a joint fit. We use logarithmic
IO coordinates, additive ambient offsets and untransformed history fractions,
with analytic mean derivatives, 24 starts, seed 532 and at most 1,200 function
evaluations per start. Appendix~\ref{sec:nuisance-assumptions} derives exact
ambiguities that can remain when these quantities are free.

Another 96 joint fits hold ambient offsets at
$-1,-0.5,-0.25,0.25,0.5,1\,{}^\circ\mathrm C$ for each delayed model and
both prescribed histories. For these fixed perturbations, the same offset
and history assumption are applied to each held-out record. Fitted
record-specific nuisance quantities are not assigned automatically to an unseen
record; those fits have no corresponding validation correction.

The remaining 49 joint fits compare input reconstruction and discretization:
32 exact left/right-hold fits across eight delayed models and two startups,
14 right-hold backward-Euler fits across the seven finite chains and two
startups, and three original-model variants (exact left hold, exact right hold,
and right-hold backward Euler). They use the same 24-start budget of 1,200
function evaluations. We retain active bounds, numerical ranks, all attempts
and any failure to recover the nested fixed-assumption objective. Reduced
calibration SSE after adding nuisance parameters supplies no independent
validation evidence.

\subsection{Local uncertainty and estimability}
For residual vector $\bm e$ and log-parameter vector $\bm\eta$, the sensitivity
matrix $J=\partial\bm e/\partial\bm\eta$ defines the conditional information
matrix $F=J^\top J/\sigma^2$ under independent and identically distributed (IID)
Gaussian errors of variance
$\sigma^2$. We computed the singular value decomposition of $J$ directly and
reported the eigenvalues and directions of $F$. Numerical rank uses a relative
singular-value cutoff of $10^{-8}$. This threshold is a numerical convention;
the explicit equivalence classes in Appendix~\ref{sec:structural} establish the
structural conclusions independently.

We reported local uncertainty under a fixed hypothetical sensor standard
deviation of $0.25\,{}^\circ\mathrm C$ and under a common variance estimated
from each calibration residual vector. The fixed value is a sensitivity scenario,
not a calibration established by these experiments. Rank-deficient full parameter
covariances are undefined. A diagnostic pseudoinverse is retained separately and
is not used as a confidence interval for a null direction. We also computed
heteroskedasticity and autocorrelation consistent (HAC) sandwich estimates at lag
bandwidths of 10, 30, and 60 samples, together with residual autocorrelation
functions (ACFs) and Ljung--Box statistics. The sandwich uses Bartlett lag weights
\citep{newey-1987-simple}. Its asymptotic justification requires zero-mean score
moments, bounded moments, and suitable decay of dependence. Fixed bandwidths
in these finite records do not establish those assumptions or nominal coverage.
These calculations assess sensitivity to
serial dependence; they do not correct a biased mean model.

\subsection{A working correlated-residual likelihood}
We also fitted the 48 IO calibration cases with stationary Gaussian
Ornstein--Uhlenbeck (OU) residual
covariance $R_{ij}=\sigma^2\exp(-|t_i-t_j|/\ell)$ within each experiment,
independently across experiments. Actual retained observation times are used.
For residuals $e_i$, set $w_1=e_1$ and
\[
 \phi_i=e^{-(t_i-t_{i-1})/\ell},\qquad
 w_i=\frac{e_i-\phi_i e_{i-1}}{\sqrt{1-\phi_i^2}}.
\]
With $q=\sum_j\|\bm w_j\|^2$ over records and $N$ retained observations,
profiling the common variance gives $\hat\sigma^2=q/N$ and negative
log-likelihood
\[
 \mathcal L=\frac12\left[
 N\{\log(2\pi)+1+\log(q/N)\}
 +\sum_j\sum_{i=2}^{N_j}\log(1-\phi_{ij}^2)\right].
\]
The determinant term is retained while optimizing the three log mean parameters
and $\log\ell$, with $\ell\in[0.05,3000]$ s, 24 starts and the same seed.
All solver attempts and boundary diagnostics are retained. We compare
conditional mean information at fixed fitted noise with observed curvature
that profiles variance and includes mean--length coupling. Two Hessian steps,
$10^{-4}$ and $5\times10^{-5}$ in log coordinates, check numerical sensitivity.
A fitted OU length or improved likelihood does not establish a residual
mechanism; raw and filtered correlation diagnostics assess this working model.

\subsection{Likelihood profiles and resampling}
For each native parameter we fixed its value and reoptimized all nuisance
parameters, following the profile construction of \citet{raue-2009-structural}.
The exploration combined a global grid with local refinement, targeting 61 points
and four nuisance starts at each point. The exact count after merging grids is
recorded per profile. Additional points examined the exchanged
pole mode of the one-stage cold-start model. We retained disconnected accepted
components and endpoints at the exploration bounds. An endpoint at such a bound
does not establish an unbounded mathematical interval.

For $N$ residuals and sum of squared errors $S(\bm\eta)$, the likelihood-ratio
statistics were $[S(\bm\eta)-S(\hat{\bm\eta})]/\sigma^2$ for fixed variance
and $N\log[S(\bm\eta)/S(\hat{\bm\eta})]$ for a common profiled variance.
The nominal 95\% reference threshold was the 0.95 quantile of a chi-squared
distribution with one degree of freedom. These thresholds assume an appropriate
independent Gaussian likelihood and regular identification; they are conditional
diagnostics when those assumptions fail.

After comparing first-pass sets with the exact physical equivalence family,
we refined a declared selection of observed crossings in a separate phase.
The selection covers the redundant original physical coordinates, cold
one-stage delay profiles in both coordinates, and the representative joint
cold higher-order delay profiles. Feasible exact-equivalence witnesses are
included before crossing refinement. Bisection stops at a likelihood-ratio
residual of $10^{-3}$ or a log-parameter bracket width of $10^{-7}$, with at
most 26 steps and four nuisance starts per new point. The original grids are
preserved. A separate audit checks forward objectives at new points and whole
compatible-family inclusion inside the fitting box; it cannot establish the
absence of other modes.

The resampling design specified 400 refits for each of four circular residual
block lengths (1, 10, 30, and 60 samples), plus 400 independent Gaussian
simulations at the fixed $0.25\,{}^\circ\mathrm C$ standard deviation, for each
input--output calibration case. Residuals were centered and resampled within
each experiment; first temperatures, ambient estimates, and heater calibration
were fixed. Each draw was transformed exactly into physical coordinates.
Structural null directions were marked undefined rather than assigned an interval
from arbitrary optimizer selections. These resampling distributions remain
conditional on the mean model and on the assumed within-experiment stationarity.
The fixed-length circular scheme retains equal marginal sampling weights at
record edges. It differs from the geometric block lengths of the stationary
bootstrap; \citet[physical PDF p.~3, Remark 2]{politis-1994-stationary} describes
the fixed-length circular relative. Neither that finite-sample mean property nor
the paper's stationary-bootstrap theory establishes coverage for our nonlinear
fits with model discrepancy. We assess selected interval procedures separately
using independent synthetic datasets.

The independent coverage design prespecified 300 synthetic datasets for each
of four joint cold generating models (one, two and twelve stages, plus pure
delay) and three true error processes: IID Gaussian, stationary OU with 10 s
length, and stationary OU with 60 s length. Every process has marginal standard
deviation $0.25\,{}^\circ\mathrm C$. We used the recorded input and observation
times, fixed the conditioned initial temperatures exactly, and refitted each
synthetic dataset before constructing local and profile intervals. The
comparison includes known-variance and estimated-variance IID log-Wald
intervals, three HAC bandwidths, and both delay-profile variance conventions.
Physical intervals use the exact parameter transformation.

All 3,600 requested datasets remain in the coverage denominator. We report
coverage among defined intervals, its Wilson binomial interval, and worst/best
unconditional bounds that count failed or undefined outcomes as uncovered or
covered, respectively. Failed worker chunks and missing result files are
distinguished: a completed chunk with a missing file fails the integrity audit.
Disconnected profile widths sum the accepted components; envelope widths also
include excluded gaps. A secondary three-point quadratic interpolation checks
whether finite-grid interpolation changes coverage classification at the known
true delay. It does not replace nuisance reoptimization or certify endpoint
accuracy. The study conditions on these generating points and does not estimate
coverage under real between-record model discrepancy.
Bootstrap intervals remain real-data sensitivity comparisons; their
repeated-sampling coverage is unmeasured by this experiment.

\subsection{Conditional prediction uncertainty}
For target mean sensitivity $\bm j(t)$, local mean variance is
$\bm j(t)C\bm j(t)^\top$. We use covariance on the identifiable quotient with
residual degrees of freedom $N-\operatorname{rank}(J)$, checking that the target
annihilates each structural null direction. Known global ambiguities require
an additional check: an equilibrated constant-input calibration leaves delay
free, so we withhold bands for a varying-input target even if its local delay
derivative happens to vanish at the selected parameter. The original model's
zero-state physical equivalence preserves predictions globally and permits
finite prediction uncertainty on its quotient.
The eleven rank-deficient baseline cases were the three original-model gauge
cases and eight equilibrated constant-step cases. The synthetic repeated-pole
example is excluded from this regular-quotient covariance interpretation.

Local pointwise normal mean bands are compared with percentile mean bands from
saved joint-fit bootstrap trajectories. An additional independent future error
law is required for observation prediction bands. We add its plug-in marginal
variance under IID or working OU assumptions; HAC supplies mean bands without
inventing a future error distribution. For Gaussian parametric bootstrap draws,
future-observation quantiles are computed from the declared equal-weight mixture
of 400 Gaussian trajectory components. This is a numerical mixture calculation,
with the inferential accuracy of the bootstrap still conditional and unproven.
Circular-block draws supply mean bands only. The first observation is excluded
from all bands because it conditions the state. Existing-record pointwise
containment is descriptive, without a repeated-experiment or simultaneous
coverage interpretation.

\subsection{Prospective experiment design}
The prespecified input-design horizon is 900 s with 1 s observations and 19
command knots at 50 s intervals, joined linearly. Command amplitudes lie between
0 and 100\%, their time average is 50\%, and the slew limit is one percentage
point per second. Ambient and initial measured temperature are fixed at
$25\,{}^\circ\mathrm C$. We constrain the predicted mean temperature below
$85\,{}^\circ\mathrm C$ for all sixteen joint fitted mean/startup combinations,
regardless of the model used for the utility. These are assumed simulation
constraints, without a hardware safety certification or uncertainty margin.
An independent 0.25 s forward simulation also checks a 300 s shutdown tail,
ramping the last command to zero at the same slew limit.
The first design trial met its 1 s constraints but exceeded the assumed limit
slightly between observations in that denser check. We retained that trial
and tightened the optimization bound by $0.01\,{}^\circ\mathrm C$ for both
parameter and discrimination designs. This numerical guard leaves the reported
acceptance limit at $85\,{}^\circ\mathrm C$; every selected design must still
pass the independent dense shutdown-tail check. The guard itself supplies no
continuous-time bound or hardware guarantee.

For a future sensitivity matrix $J(u)$ and assumed covariance $R$, the local
information is $F(u)=J(u)^\top R^{-1}J(u)$, optionally augmented by information
from the two historical calibration experiments under the same model and
covariance. This added information is not a Bayesian distribution over
parameters. The four utilities are $\log\det F$, $\log\lambda_{\min}(F)$,
$-\log[(F^{-1})_{\theta\theta}]$, and
$-\log\operatorname{tr}(A F^{-1})$. Here $A$ is the scenario-specific average
thermal-mean sensitivity outer product for a future constant 50\% reference.
The variance target concerns mean prediction, without future measurement noise.
The third utility targets log-delay variance.
Null-space compatibility is checked before a target-specific pseudoinverse is
used. Parameter metrics are natural logarithms of $(K,\tau,\theta)$.

Nominal designs use cold one-stage, two-stage, twelve-stage and pure-delay
fits with IID, OU10 s or OU60 s working correlations, each with marginal
standard deviation $0.25\,{}^\circ\mathrm C$. Ensemble designs use all 48
model/startup/covariance scenarios. We aggregate scenario-relative log utilities
by their mean or minimum. Four utilities and two historical-information settings
give 112 planned designs. Six common command starts include constant, sine,
multisine, random binary and alternating knot patterns, each satisfying the
amplitude, dose and slew rules. Constant input supplies the normalization
reference where its utility is finite; otherwise we use the multisine.
Multistart sequential least squares programming (SLSQP) uses
analytic input gradients and retains convergence and constraint violations;
its local solutions have no global-optimality certificate. Minimum aggregation
is nonsmooth at scenario ties.

The physical and IO log-coordinate information matrices are related by the
congruence in Appendix~\ref{sec:structural}. Because $|\det M|=1$, the
log-determinant utility is equal at corresponding points. Delay variance and
properly transformed prediction loss are also invariant. The smallest
eigenvalue depends on the chosen metric. Thus a coordinate change can alter
an E-criterion without changing the underlying prediction experiment.

We separate parameter precision from discrimination of stage count. For two
fixed future means with the same assumed covariance, define
\[
 D_{ij}^2(u)=[\mu_i(u)-\mu_j(u)]^\top R^{-1}[\mu_i(u)-\mu_j(u)].
\]
The Gaussian log-likelihood ratio has mean $\pm D_{ij}^2/2$ and variance
$D_{ij}^2$ under these two simple hypotheses. Equal-prior minimum classification
error is therefore $\Phi(-D_{ij}/2)$. This ideal calculation assumes fixed known
means and covariance; it is not a testing-power estimate after parameter
estimation or under TC-Lab discrepancy.

Twelve additional local designs aggregate relative log squared separation for
all six pairs among one-stage, two-stage, twelve-stage and pure-delay means,
at both startups and all three working covariances. The multisine supplies the
normalizing reference. A bounded 72-case diagnostic then lets each alternative
model refit its three parameters against a fixed one-stage generating mean,
using 24 starts under the same future covariance. It compares constant,
multisine, and both selected discrimination designs against two-stage,
twelve-stage and pure-delay alternatives. This directional subset does not
cover every generating model or parameter point. The one-stage order was chosen
as a representative because it has the smallest cold joint calibration error;
this choice does not establish it as the true system. A converged refit distance
is an upper bound on the globally closest alternative mean; we assign it no
simple-hypothesis error probability or chi-squared calibration.

No proposed waveform has been applied to laboratory hardware.

% End expanded input


% Expanded input: sections/results
\section{Results}
The baseline comparison completed all 99 cases. The first profile pass completed
297 of 300 jobs; a separate documented retry recovered the three failed
original-model conductance profiles. Both passes are retained. None of the
completed profiles found a minimum requiring a baseline refit under the recorded
tolerance. The bootstrap completed all 96,000 planned refits, and the working
OU likelihood comparison completed all 48 calibrations. Synthetic coverage, ambient/history and numerical sensitivity, and prediction uncertainty also completed their declared matrices and independent arithmetic audits. The parameter-design and model-discrimination matrices also completed their declared cases and independent audits.

\subsection{Order and initialization}
The startup assumption reverses the order ranking for the sine and joint
calibrations (Fig.~\ref{fig:baseline-order}). For these two calibration sets,
cold delay states give the lowest error at one stage among the compared
formulations. Equilibrated states favor
the pure-delay limit. The step experiment with equilibrated states gives the
same error for every order, as expected from the delay's absence from its
response. These comparisons concern calibration error under prescribed histories;
they do not establish which history occurred before recording began.
Table~\ref{tab:calibration-rmse} gives the numerical comparison, including the
step calibration, where the cold two-stage model improves on the cold one-stage
fit. Table~\ref{tab:joint-parameters} gives the joint cold parameters.

% Expanded input: tables/calibration_rmse
% Generated by render_tables.py; edit the producer, not this table.
\begin{table}[tbp]
\centering
\small
\begin{tabular}{lrrrr}
\toprule
Model & Sine, cold & Step, cold & Joint, cold & Joint, equilibrated \\
\midrule
$n=1$ & 0.2567 & 0.1573 & 0.4115 & 0.7936 \\
$n=2$ & 0.3470 & 0.1529 & 0.4630 & 0.7585 \\
$n=3$ & 0.3810 & 0.1573 & 0.4841 & 0.7486 \\
$n=4$ & 0.3985 & 0.1610 & 0.4953 & 0.7441 \\
$n=6$ & 0.4165 & 0.1659 & 0.5070 & 0.7398 \\
$n=8$ & 0.4256 & 0.1689 & 0.5130 & 0.7378 \\
$n=12$ & 0.4348 & 0.1724 & 0.5192 & 0.7359 \\
Pure delay & 0.4541 & 0.1812 & 0.5323 & 0.7322 \\
\bottomrule
\end{tabular}
\caption{Calibration RMSE in degrees Celsius, excluding the initial temperature. Each entry comes from its own fitted calibration/startup combination. Equilibrated means the prescribed delay states equal the initial command, with the sensor temperature fixed by the first observation.}
\label{tab:calibration-rmse}
\end{table}

% End expanded input


% Expanded input: tables/joint_parameters
% Generated by render_tables.py; edit the producer, not this table.
\begin{table}[tbp]
\centering
\small
\begin{tabular}{lrrrrr}
\toprule
Model & $K$ & $\tau$ (s) & $\theta$ (s) & $\Ua$ (W/K) & $\Cp$ (J/K) \\
\midrule
$n=1$ & 0.7744 & 150.697 & 16.555 & 0.04132 & 6.227 \\
$n=2$ & 0.7762 & 154.197 & 15.403 & 0.04123 & 6.357 \\
$n=3$ & 0.7768 & 155.291 & 15.074 & 0.04119 & 6.397 \\
$n=4$ & 0.7771 & 155.826 & 14.917 & 0.04118 & 6.417 \\
$n=6$ & 0.7774 & 156.353 & 14.764 & 0.04116 & 6.436 \\
$n=8$ & 0.7776 & 156.614 & 14.689 & 0.04115 & 6.445 \\
$n=12$ & 0.7777 & 156.874 & 14.614 & 0.04115 & 6.455 \\
Pure delay & 0.7780 & 157.391 & 14.466 & 0.04113 & 6.474 \\
\bottomrule
\end{tabular}
\caption{Joint cold-start estimates. Gain units are degrees Celsius per percentage-point command. Physical columns transform the IO estimates using fixed $c=0.032$ W per percentage point; they are effective-model quantities. The one-stage row uses the representative $\tau\geq\theta$.}
\label{tab:joint-parameters}
\end{table}

% End expanded input


\begin{figure}[p]
 \centering
 \includegraphics[width=.5\linewidth]{baseline_order.pdf}
 \caption{Initialization changes the fitted order ranking. Root mean squared
 temperature error (RMSE) for sine, step, and joint calibration is plotted for
 input--output parameter fits. Cold means $z_i(0)=0$; equilibrated means
 $z_i(0)=u(0)$. The infinity label denotes the separately fitted pure-delay model,
 with the corresponding input prehistory. Lines connect categorical orders as
 guides. Every error excludes the first temperature used to initialize the
 model. No uncertainty intervals are implied by the markers.}
 \label{fig:baseline-order}
\end{figure}

\subsection{Temperature traces and residuals}
The jointly calibrated cold models give closely spaced temperature predictions,
while their residual traces show differences near startup and the sine peaks
(Fig.~\ref{fig:baseline-trajectories}). Residual structure persists over many
successive observations. A visual match on the full temperature scale therefore
does not establish independent sensor noise or reliable uncertainty coverage.

\begin{figure}[p]
 \centering
 \includegraphics[width=\linewidth]{baseline_trajectories.pdf}
 \caption{The residual scale exposes differences hidden in the temperature
 overlays. Columns show the sine and step experiments; rows show heater input
 $u$, measured temperature $T$ with jointly fitted cold-start mean predictions,
 and residuals defined as prediction minus observation. All recorded calibration
 samples are plotted. The thermal initial condition and assumed ambient equal
 the first observation, and the delay prehistory is zero. The four curves use
 independent joint fits for their respective formulations. These are mean
 predictions; uncertainty bands are not shown.}
 \label{fig:baseline-trajectories}
\end{figure}

\subsection{Transfer to supplied validation files}
The four representative joint cold formulations give similar errors on the
eleven duplicate-excluded validation files (Fig.~\ref{fig:validation-errors}).
Differences among files are much larger than differences among these stage
counts. Most pooled squared error is accounted for algebraically by the
per-file mean residuals (Table~\ref{tab:validation-summary}). These offsets
are retrospective diagnostics, without an identified physical cause or a
prospective correction. Ambient, command history, heater conversion, thermal
coefficients and omitted dynamics remain possible explanations. The supplied
file names alone do not distinguish them.

% Expanded input: tables/validation_summary
% Generated by plot_validation.py; see validation_figure_manifest.json.
\begin{table}[tbp]
\centering\small
\begin{tabular}{lrrr}
\toprule
Model & Calibration RMSE & Held-out RMSE & File-mean share (\%) \\
\midrule
$n=1$ & 0.4115 & 3.3588 & 84.8 \\
$n=2$ & 0.4630 & 3.3608 & 83.4 \\
$n=12$ & 0.5192 & 3.3883 & 81.1 \\
Pure delay & 0.5323 & 3.4065 & 80.1 \\
\bottomrule
\end{tabular}
\caption{Joint cold baseline fits evaluated without refitting on eleven duplicate-excluded validation files. RMSE values are in degrees Celsius; the held-out value pools all retained observations rather than averaging file RMSEs. File-mean share is $100\sum_j N_j\bar e_j^2/\sum_j\sum_i e_{ij}^2$, the retrospective algebraic contribution of each file's mean residual to pooled held-out error. It is not a causal attribution or a prospective bias correction.}
\label{tab:validation-summary}
\end{table}

% End expanded input


% Expanded input: tables/validation_file_map
% Generated by plot_validation.py; see validation_figure_manifest.json.
\begin{table}[tbp]
\centering\small
\begin{tabular}{llrrr}
\toprule
Code & Supplied filename suffix & Duration (s) & $T_0$ (${}^\circ$C) & $N$ \\
\midrule
V01 & \path{env_2_sin_2_50_run_1.csv} & 900.0 & 20.91 & 899 \\
V02 & \path{env_2_sin_2_50_run_2.csv} & 900.0 & 20.26 & 900 \\
V03 & \path{env_2_sin_2_50_run_3.csv} & 900.0 & 20.26 & 900 \\
V04 & \path{env_2_sin_50_50_run_1.csv} & 900.0 & 19.94 & 900 \\
V05 & \path{env_2_sin_50_50_run_2.csv} & 900.0 & 20.58 & 898 \\
V06 & \path{env_2_sin_50_50_run_3.csv} & 900.0 & 20.26 & 898 \\
V07 & \path{env_2_sin_5_50_run_2.csv} & 900.0 & 23.48 & 900 \\
V08 & \path{env_2_sin_5_50_run_3.csv} & 900.0 & 24.13 & 900 \\
V09 & \path{env_2_step_50_run_2.csv} & 900.0 & 24.13 & 900 \\
V10 & \path{env_2_step_50_run_3.csv} & 900.0 & 23.81 & 900 \\
V11 & \path{sin_5_50.csv} & 1000.0 & 24.77 & 989 \\
\bottomrule
\end{tabular}
\caption{Validation-file codes. Restore the common prefix \texttt{validation\_experiment\_} to each supplied filename. $N$ excludes the initial temperature, which also supplies the assumed ambient for that file; it is not an independent ambient measurement. File labels do not establish independent replication or a change in physical environment.}
\label{tab:validation-files}
\end{table}

% End expanded input


\begin{figure}[p]
 \centering
 \includegraphics[width=\linewidth]{validation_errors.pdf}
 \caption{Transfer error varies more among supplied files than among the
 representative cold delay formulations. Panels show RMSE and signed mean
 residual, defined as prediction minus observation, excluding each initial
 temperature. Parameters come from the two joint calibration records and
 are not refitted on these files. Codes are defined in
 Table~\ref{tab:validation-files}; lines join categorical file codes as guides.
 Each prediction uses the file's first temperature as both its thermal initial
 condition and assumed ambient. Markers do not imply independent replication,
 and no uncertainty interval is shown.}
 \label{fig:validation-errors}
\end{figure}

\subsection{Coordinate conditioning and physical ambiguity}
The independently fitted physical and input--output formulations reach sums of
squares that agree to numerical precision in the joint cold cases
(Table~\ref{tab:coordinate-metrics}). Reexpressing the same sensitivity matrix
changes its condition number. Moment coordinates do not improve this particular
metric at these fitted points. These differences concern the coordinate metric;
they do not distinguish observationally equivalent model explanations.

% Expanded input: tables/coordinate_metrics
% Generated by render_tables.py; edit the producer, not this table.
\begin{table}[tbp]
\centering
\small
\begin{tabular}{lrrrr}
\toprule
Model & IO logs & Physical logs & Moment & $|\Delta S|$ \\
\midrule
$n=1$ & 23.1 & 19.38 & 27.76 & $6.3\times10^{-12}$ \\
$n=2$ & 21.04 & 17.56 & 24.82 & $1.5\times10^{-11}$ \\
$n=3$ & 20.61 & 17.19 & 24.15 & $1.3\times10^{-12}$ \\
$n=4$ & 20.44 & 17.04 & 23.86 & $4.3\times10^{-11}$ \\
$n=6$ & 20.29 & 16.9 & 23.6 & $2.1\times10^{-11}$ \\
$n=8$ & 20.23 & 16.84 & 23.48 & $3.2\times10^{-12}$ \\
$n=12$ & 20.16 & 16.78 & 23.36 & $3.4\times10^{-13}$ \\
Pure delay & 20 & 16.64 & 23.09 & $8.0\times10^{-11}$ \\
\bottomrule
\end{tabular}
\caption{Coordinate metrics for joint cold fits. The first three columns are sensitivity condition numbers at the same IO model point, transforming the information matrix without refitting. Moment coordinates are $(\log K,\log m,\ell)$ with $m=\tau+\theta$ and $\ell=\log[\theta/\tau]$. The last column compares independently fitted IO and physical sums of squares, in squared degrees Celsius. Conditioning changes do not imply additional information.}
\label{tab:coordinate-metrics}
\end{table}

% End expanded input

Table~\ref{tab:coordinate-correlations} shows the corresponding local parameter
correlations. Transforming the capacity combines gain and time-constant
uncertainty, so the physical-coordinate correlations need not resemble the IO
correlations even though the predicted thermal response is unchanged.

% Expanded input: tables/coordinate_correlations
% Generated by render_structural_tables.py; edit the producer, not this table.
\begin{table}[tbp]
\centering
\small
\begin{tabular}{lrrr}
\toprule
Model and log coordinates & $\rho_{12}$ & $\rho_{13}$ & $\rho_{23}$ \\
\midrule
$n=1$, IO & 0.682 & -0.366 & -0.696 \\
$n=1$, Physical & -0.535 & 0.366 & -0.711 \\
$n=2$, IO & 0.678 & -0.312 & -0.602 \\
$n=2$, Physical & -0.511 & 0.312 & -0.616 \\
$n=12$, IO & 0.676 & -0.269 & -0.523 \\
$n=12$, Physical & -0.494 & 0.269 & -0.535 \\
Pure delay, IO & 0.675 & -0.260 & -0.506 \\
Pure delay, Physical & -0.491 & 0.260 & -0.518 \\
\bottomrule
\end{tabular}
\caption{Local IID correlations at the same joint cold model point. IO coordinates are $(\log K,\log\tau,\log\theta)$; physical coordinates are $(\log\Ua,\log\Cp,\log\theta)$. The physical covariance is transformed without refitting. Correlation changes reflect the chosen quantities, not extra information. One-stage entries describe the local canonical pole mode, excluding its exchanged global label.}
\label{tab:coordinate-correlations}
\end{table}

% End expanded input


The positive information spectra are similar across the reduced cold models,
while their native local dimensions differ from the original model and the
equilibrated step case (Fig.~\ref{fig:information-spectra}). The weakest local
direction is dominated by relative delay changes at these joint cold points,
with smaller compensating time-constant changes. Its components change after
reexpressing the same model in physical coordinates
(Fig.~\ref{fig:weak-directions}). Neither local eigenanalysis detects the
one-stage pole-label exchange by itself.

\begin{figure}[p]
 \centering
 \includegraphics[width=.85\linewidth]{information_spectra.pdf}
 \caption{Positive information eigenvalues and structural null directions must
 be counted separately. Values above the relative sensitivity-rank cutoff are
 normalized by each case's largest eigenvalue. Legend parentheses give local
 identifiable rank divided by native parameter dimension. The four reduced
 curves use joint cold fits in log input--output coordinates and overlap closely.
 The original joint model uses four physical logs; the step comparison uses
 two-stage equilibrated delay states. Their null eigenvalues are omitted rather
 than assigned a plotting floor. These coordinate-dependent spectra are not an
 invariant ranking of information magnitude between formulations. Lines join
 ordered eigenvalues as guides.}
 \label{fig:information-spectra}
\end{figure}

\begin{figure}[p]
 \centering
 \includegraphics[width=\linewidth]{weak_directions.pdf}
 \caption{Weakest local information directions depend on the coordinate metric.
 Columns give joint cold fits for the indicated stage count, with infinity
 denoting pure delay. Entries are components of unit eigenvectors for the
 smallest information eigenvalue. The largest absolute component is assigned
 positive sign. Physical-coordinate directions are re-diagonalized after the
 exact transformation at the same fitted point; they are not additional
 observations. The displayed directions concern one local pole representative
 and do not remove global parameter ambiguities.}
 \label{fig:weak-directions}
\end{figure}

The original model has a stronger ambiguity. For each calibration, the saved
structural calculation sampled 403 positive physical parameter combinations with
the same transfer function. Their largest prediction difference was below
$1.8\times10^{-12}\,{}^\circ\mathrm C$. Conductance $\Ua$ remained fixed while
$\Ub$, $\CpH$ and $\CpS$ varied. Appendix~\ref{sec:structural} gives the
analytic family and its positivity limits; sampled ranges are not confidence
intervals.

Figure~\ref{fig:physical-family} shows the joint family and the forward numerical
agreement along it. A selected optimizer point is one compatible member.
By contrast, the family's limiting capacities and maximum coupling are functions
of identifiable transfer coefficients. Table~\ref{tab:family-limits} and
Fig.~\ref{fig:family-boundaries} propagate saved one-stage bootstrap draws to
these boundary functions. Larger residual blocks widen their conditional
percentile intervals. This quantifies uncertainty in the family limits without
assigning an arbitrary member the status of a physical estimate.

% Expanded input: tables/family_limits
% Generated by render_structural_tables.py; edit the producer, not this table.
\begin{table}[tbp]
\centering
\small
\begin{tabular}{lrrr}
\toprule
Boundary function & Fitted value & IID simulation & Circular block 60 \\
\midrule
$\inf\CpH$ (J/K) & 0.684 & $[0.6792,0.6896]$ & $[0.6448,0.7276]$ \\
$\sup\CpH$ (J/K) & 6.227 & $[6.214,6.238]$ & $[6.144,6.315]$ \\
$\sup\CpS$ (J/K) & 2.783 & $[2.765,2.799]$ & $[2.647,2.91]$ \\
$\sup\Ub$ (W/K) & 0.07451 & $[0.07359,0.07528]$ & $[0.06801,0.08156]$ \\
\bottomrule
\end{tabular}
\caption{Identifiable boundary functions for the full strictly positive original-model family, propagated from the joint cold one-stage transfer fit. Intervals are conditional 95 percent percentiles from 400 saved draws per method, before numerical physical-parameter boxes. The heater-capacity infimum and supremum are open family endpoints; the sensor-capacity and coupling maxima are attained. These are not marginal intervals for an arbitrary physical parameter.}
\label{tab:family-limits}
\end{table}

% End expanded input


\begin{figure}[p]
 \centering
 \includegraphics[width=\linewidth]{physical_family.pdf}
 \caption{Different original physical parameters give the same measured
 temperature response. The joint cold transfer function generates 403 positive
 members parameterized by $r=\CpS/\Ub$. Panels (a)--(c) show the heater capacity,
 sensor capacity and coupling conductance. Orange dots mark the member selected
 by the original optimizer. Panel (d) shows maximum absolute prediction
 differences over both calibration records, exposing floating-point agreement
 rather than measurement precision. Of the displayed members, 401 fall inside
 the original numerical fitting box. Strictly positive family endpoints are
 approached but not attained; the curves are not confidence regions.}
 \label{fig:physical-family}
\end{figure}

\begin{figure}[p]
 \centering
 \includegraphics[width=\linewidth]{family_boundary_intervals.pdf}
 \caption{Uncertainty in identifiable physical-family boundaries can be
 propagated without selecting a physical family member. Panels show the
 heater-capacity infimum and supremum, sensor-capacity supremum and coupling
 supremum. Each bar is a conditional 95\% percentile interval from 400 saved
 joint cold one-stage bootstrap draws, expressed as percent change from its
 fitted boundary value; dots mark medians. The strictly positive family
 precedes rectangular numerical fit boxes. These are intervals for boundary
 functions, excluding heater-calibration and model-discrepancy uncertainty;
 they are not marginal intervals for $\CpH$, $\CpS$ or $\Ub$.}
 \label{fig:family-boundaries}
\end{figure}

The analytic family also exposed an interpolation failure in the first-pass
profiles. All nine redundant physical-parameter grid sets excluded some known
equivalent members. Their sampled objectives were correct, but interpolation
across the rising edges gave misleading endpoints. A separate refinement added
verified equivalence witnesses and bisected the observed threshold crossings.
The selected catalog now includes the complete analytically compatible intervals
inside the numerical fitting box for all nine profiles. First-pass files remain
available, and the revised numerical endpoints do not make the individual
physical parameters identifiable. The same selected refinement treats the cold
one-stage delay profiles in both coordinates and all three calibration sets,
and the representative joint cold two-stage, twelve-stage and pure-delay
profiles. Its arithmetic, newly evaluated forward objectives and family inclusion
passed a separate audit. These checks do not certify unobserved modes or global
nuisance minima.

\subsection{Conditional delay uncertainty}
Table~\ref{tab:delay-uncertainty} compares delay uncertainty for representative
joint cold formulations. Increasing the residual block length widens the
percentile intervals relative to independent Gaussian simulation. The one-stage
profile has two accepted pole-label modes; canonical bootstrap summaries refer
to only the representative with $\tau\geq\theta$. A finite full-coordinate
interval is withheld where the recorded experiment leaves delay unobservable.
The table compares assumptions rather than estimating a single unconditional
physical uncertainty.

% Expanded input: tables/delay_uncertainty
% Generated by render_tables.py; edit the producer, not this table.
\begin{table}[tbp]
\centering
\small
\begin{tabular}{lcccc}
\toprule
Model & IID profile set & IID simulation & Circular block 60 & OU local SD \\
\midrule
$n=1$ & $[16.35,16.76]\cup[150.2,151.2]$ & $[16.44,16.69]$ & $[15.62,17.56]$ & 0.585 \\
$n=2$ & $[15.2,15.6]$ & $[15.3,15.51]$ & $[14.31,16.44]$ & 0.606 \\
$n=12$ & $[14.41,14.82]$ & $[14.52,14.72]$ & $[13.25,15.98]$ & 0.654 \\
Pure delay & $[14.26,14.67]$ & $[14.38,14.56]$ & $[13.34,15.68]$ & 0.317 \\
\bottomrule
\end{tabular}
\caption{Joint cold delay uncertainty in seconds. Profile sets use a profiled common IID variance and retain both one-stage pole modes. Percentile intervals use 400 Gaussian simulations at fixed $\sigma=0.25$ degrees Celsius or 400 circular residual-block refits, respectively. Bootstrap one-stage delays use the canonical pole representative. OU entries are local standard deviations with variance profiled and length included in the Hessian, not interval endpoints. These procedures make different assumptions and are not interchangeable coverage guarantees.}
\label{tab:delay-uncertainty}
\end{table}

% End expanded input


Figure~\ref{fig:profile-modes} displays the two one-stage profile modes on
separate delay windows. Figure~\ref{fig:delay-intervals} compares the conditional
intervals across residual block lengths. The apparent precision changes with
the residual assumptions even when the fitted thermal trajectories remain close.

\begin{figure}[p]
 \centering
 \includegraphics[width=\linewidth]{profile_modes.pdf}
 \caption{The one-stage cold model has two distinct delay-label modes. Each
 panel shows a separate declared window of the joint delay profile; the
 intervening parameter range is omitted. The blue solid curve uses fixed
 hypothetical sensor standard deviation $0.25\,{}^\circ\mathrm C$, and the
 orange dashed curve profiles a common IID variance. Lines connect saved
 nuisance-optimized grid values. The dotted line is the nominal 95\%
 chi-squared threshold with one degree of freedom. Negative likelihood-ratio
 roundoff is clamped to zero for display. Neither regular threshold assumptions
 nor real-data coverage follow from the profile shape.}
 \label{fig:profile-modes}
\end{figure}

\begin{figure}[p]
 \centering
 \includegraphics[width=\linewidth]{delay_intervals.pdf}
 \caption{Residual assumptions change conditional delay intervals. All panels
 use joint cold fits; the one-stage panel displays the canonical short-delay
 representative. Bootstrap bars are percentile intervals from 400 saved draws
 per method, with dots at medians. The profile bar uses the unknown-variance IID
 accepted component containing the baseline estimate. OU bars exponentiate
 local normal log-delay endpoints from the observed mean/length Hessian with
 variance profiled. Their dots are fitted OU delays. Gray vertical lines mark
 baseline least-squares delays. Residual blocks count retained samples. These
 are nominal conditional intervals under different assumptions, not a common
 unconditional physical uncertainty or simultaneous confidence statement.}
 \label{fig:delay-intervals}
\end{figure}


% Expanded input: sections/coverage_results
\subsection{Synthetic coverage exposes sensitivity to serial dependence}

% Expanded input: tables/coverage_values
% Direct fields formatted from the audited coverage summary.
\newcommand{\CoverageReplicates}{300}
\newcommand{\CoverageTotalDatasets}{3600}
\newcommand{\CoverageTwoStageOUKnownGridChanges}{68}

% End expanded input

All \CoverageTotalDatasets{} prespecified synthetic datasets finished without
failed draws. Each of the 84 delay-method comparisons returned an interval for
every requested dataset, so its method-defined and unconditional coverage rates
coincide. Table~\ref{tab:coverage-delay} gives exact counts for four procedures;
Fig.~\ref{fig:coverage-delay} includes all seven and their Monte Carlo uncertainty.

Under IID Gaussian generation, the IID local and profile procedures gave
coverage near the nominal reference. Applying those same procedures under OU
generation produced severe undercoverage, especially at the longer correlation
length. Estimating a common variance did not compensate for ignoring serial
covariance. HAC intervals improved coverage under correlated generation, with
substantial bandwidth sensitivity. Even the largest tested bandwidth fell below
the nominal reference in the OU60 comparisons. The synthetic thermal mean is
correct by construction here; these failures do not require between-record
mean discrepancy.

% Expanded input: tables/coverage_rates
% Generated by render_coverage_table.py; see coverage_table_manifest.json.
\begin{table}[tbp]
\centering\small
\begin{tabular}{llrrrr}
\toprule
Model & Error law & IID known & IID fitted & HAC60 & Profile fitted \\
\midrule
$n=1$ & IID & 278/300 & 278/300 & 273/300 & 278/300 \\
$n=1$ & OU10 & 105/300 & 103/300 & 263/300 & 102/300 \\
$n=1$ & OU60 & 73/300 & 66/300 & 235/300 & 66/300 \\
$n=2$ & IID & 281/300 & 280/300 & 272/300 & 280/300 \\
$n=2$ & OU10 & 110/300 & 110/300 & 266/300 & 108/300 \\
$n=2$ & OU60 & 61/300 & 56/300 & 249/300 & 55/300 \\
$n=12$ & IID & 282/300 & 281/300 & 276/300 & 280/300 \\
$n=12$ & OU10 & 106/300 & 104/300 & 262/300 & 104/300 \\
$n=12$ & OU60 & 59/300 & 53/300 & 248/300 & 53/300 \\
Pure delay & IID & 290/300 & 288/300 & 283/300 & 288/300 \\
Pure delay & OU10 & 106/300 & 104/300 & 273/300 & 103/300 \\
Pure delay & OU60 & 74/300 & 68/300 & 254/300 & 68/300 \\
\bottomrule
\end{tabular}
\caption{Delay-interval coverage counts for four selected procedures. Each entry is covered/defined, from \CoverageReplicates{} requested datasets at that generating point and true covariance; -- would denote no defined interval. All intervals have nominal 95\% reference levels. IID known/fitted refer to Gaussian log-Wald intervals; HAC60 uses a Bartlett bandwidth of 60 retained samples; profile fitted uses a common profiled IID variance and retains disconnected components. OU10/OU60 are true correlation lengths in seconds, with fixed marginal standard deviation. Full method-specific failure bounds and grid diagnostics are retained in the coverage CSV; Wilson Monte Carlo intervals and all seven methods appear in Fig.~\ref{fig:coverage-delay}.}
\label{tab:coverage-delay}
\end{table}

% End expanded input


The profile comparison retains a numerical limitation. Three-point quadratic
interpolation changed \CoverageTwoStageOUKnownGridChanges{} of
\CoverageReplicates{} coverage classifications for the two-stage known-variance
profile under OU60 generation. This is a sensitivity check on the saved grid,
without new nuisance optimization or an endpoint-accuracy certificate. Both
primary and secondary classifications remain conditional diagnostics of an IID
likelihood applied to correlated errors. The full CSV retains this diagnostic
for every profile case; interpreting the known-variance profile rates as exact
continuous-profile coverage would overstate their numerical verification.

Figure~\ref{fig:coverage-width} compares median interval widths. The one-stage
profile width sums its two accepted pole-label components rather than measuring
the gap between them. A wider accepted set can therefore concern a different
global parameter representation from a local interval around one pole label.
These simulations fix four generating points, the recorded inputs/times and the
initial-temperature, ambient and heater-calibration assumptions. They do not
measure coverage of the real-data block-bootstrap intervals or uncertainty from
the observed between-record discrepancy.

\begin{figure}[p]
 \centering
 \includegraphics[width=\linewidth]{coverage_delay_coverage.pdf}
 \caption{Serial dependence reduces coverage of IID-based delay intervals in
 controlled simulation. Each point is covered/defined over
 \CoverageReplicates{} requested datasets for one model and true error process;
 bars are pointwise 95\% Wilson intervals for Monte Carlo sampling uncertainty.
 All displayed procedures returned intervals for every dataset. The dashed line
 is nominal 0.95 coverage. Known/fitted denote fixed or profiled IID variance;
 local IID intervals are Gaussian log-Wald intervals and profiles retain all
 accepted components. HAC labels give Bartlett bandwidths in retained samples.
 All true errors have the same marginal standard deviation; rows change their
 correlation law. The thermal mean and startup are correctly specified in each
 generating scenario. Profile sets retain the declared finite-grid limitation.
 These are generating-point results, without a universal or real-discrepancy
 coverage guarantee.}
 \label{fig:coverage-delay}
\end{figure}

\begin{figure}[p]
 \centering
 \includegraphics[width=\linewidth]{coverage_delay_width.pdf}
 \caption{Median delay-interval widths depend on the uncertainty procedure and
 true serial dependence. Models, methods and simulation denominators match
 Fig.~\ref{fig:coverage-delay}. All panels share the same width scale. Profile
 widths sum accepted component lengths, excluding intervening gaps; the
 one-stage profile includes both pole-label modes. These width summaries do
 not establish nominal coverage. Exact counts, width distributions, exploration
 flags and interpolation sensitivity remain in the saved coverage summary.}
 \label{fig:coverage-width}
\end{figure}

% End expanded input


\subsection{Nonregular curvature and the delay-chain limit}
Two deterministic checks distinguish local singularity from exact physical
ambiguity (Fig.~\ref{fig:nonregular-limits}). Near a repeated pole, the squared
temperature distance grows quartically along the fixed-mean perturbation,
consistent with the canceled first derivative in Appendix~\ref{sec:structural}.
The rank-deficient quadratic approximation misses this higher-order distinction.
Separately, the cold constant-step onset excludes a uniform $O(n^{-1})$
approximation error bound. At $t=\theta$, the pure-delay response has not risen,
whereas the finite-chain rise decays as $n^{-1/2}$ with the derived coefficient.
These are fixed-parameter model implications, rather than comparisons of
refitted TC-Lab errors.

\begin{figure}[p]
 \centering
 \includegraphics[width=\linewidth]{nonregular_limits.pdf}
 \caption{Two deterministic limits challenge routine local approximations.
 (a) Squared prediction distance over the recorded sine command for
 $K=0.78$, $\tau=10+\delta$ and $\theta=10n-\delta$, with cold states.
 The dotted quartic guide is anchored at the smallest positive two-stage
 perturbation. (b) Exact cold constant-step temperature rise at $t=\theta$
 for $K=0.78$, $\tau=150$ s, $\theta=18$ s and $u=50$ percentage points.
 The reference coefficient is $A=Ku\theta/(\tau\sqrt{2\pi})$; the pure-delay
 rise is zero at that instant. The calculations use separate parameter and
 input settings and do not constitute fitted order selection or uncertainty
 coverage.}
 \label{fig:nonregular-limits}
\end{figure}

\subsection{The working correlated-noise model remains inadequate}
The independent innovations audit reproduced every saved OU likelihood within
$2.6\times10^{-12}$. This validates likelihood arithmetic, not the residual
model. Filtered innovations retain strong negative lag-one correlation
(Table~\ref{tab:ou-diagnostics}). The fitted length therefore cannot be treated
as a measured thermal disturbance time scale. Pooled held-out prediction error
also remains much larger than calibration error. Allowing
serial covariance changes conditional uncertainty without establishing that
the thermal mean transfers between records.

% Expanded input: tables/ou_diagnostics
% Generated by render_tables.py; edit the producer, not this table.
\begin{table}[tbp]
\centering
\small
\begin{tabular}{lrrrr}
\toprule
Model & Length (s) & $\hat\sigma$ (${}^\circ$C) & ACF1, sine & ACF1, step \\
\midrule
$n=1$ & 17.46 & 0.4088 & -0.5333 & -0.492 \\
$n=2$ & 22.14 & 0.459 & -0.5291 & -0.4935 \\
$n=12$ & 27.91 & 0.5151 & -0.5253 & -0.4877 \\
Pure delay & 29.33 & 0.5289 & -0.518 & -0.4811 \\
\bottomrule
\end{tabular}
\caption{Working OU fits to the two joint cold records. The scale is the maximum-likelihood stationary residual standard deviation. ACF1 is the centered sample lag-one correlation of filtered innovations. Strong remaining negative correlation challenges this one-component residual covariance; a high fitted likelihood does not validate it.}
\label{tab:ou-diagnostics}
\end{table}

% End expanded input

Figure~\ref{fig:ou-correlations} compares raw residuals and filtered innovations
within each OU fit. For reference, filtering an IID component at regular
sampling gives interior lag-one correlation $-\phi/(1+\phi^2)$, approaching
$-1/2$ for $\phi$ near one (Appendix~\ref{sec:nuisance-assumptions}). This
provides one explanation to investigate, without identifying a unique error
mechanism from the observed negative correlation.

\begin{figure}[p]
 \centering
 \includegraphics[width=\linewidth]{ou_correlations.pdf}
 \caption{A one-component OU covariance removes much low-frequency residual
 correlation while leaving strongly negative first-lag innovations. Columns
 show calibration records; rows show centered sample autocorrelations before
 and after the stationary OU filter. Each curve uses its own jointly fitted
 cold thermal mean and OU length in both rows, isolating the filter comparison
 within that fit. Lag zero is omitted. These are empirical residual
 diagnostics, without nominal whiteness-test bands or a unique physical
 interpretation of the residual process.}
 \label{fig:ou-correlations}
\end{figure}


% Expanded input: sections/followup_results
\subsection{Ambient and history assumptions shift the fitted delay}
Allowing record-specific ambient offsets reduces joint calibration error more
than allowing only an initial history fraction in the representative models
(Fig.~\ref{fig:nuisance-sensitivity}). The fitted delay also changes, although
the cold one-stage formulation retains the lowest calibration error among these
four representatives. With both nuisance terms free, every displayed step-record
ambient offset reaches the imposed upper bound
(Table~\ref{tab:nuisance-parameters}). The lower residual error therefore does
not identify an independently measured ambient temperature. Extending that bound
until the fit looks satisfactory would change the sensitivity question rather
than resolve the missing measurement.

% Expanded input: tables/nuisance_parameters
% Generated by report_followups.py; scalar sources in followup_plot_data.json.
\begin{table}[tbp]
\centering\small
\begin{tabular}{lrrrrrr}
\toprule
Model & \shortstack{$K$\\(${}^\circ$C/\%)} & $\theta$ (s) & \shortstack{$\Delta T_{a,\sin}$\\(${}^\circ$C)} & \shortstack{$\Delta T_{a,\rm step}$\\(${}^\circ$C)} & $\gamma_{\sin}$ & $\gamma_{\rm step}$ \\
\midrule
$n=1$ & 0.7390 & 17.159 & 1.339 & 2.000 & 0.000 & 0.004 \\
$n=2$ & 0.7406 & 16.350 & 1.336 & 2.000 & 0.088 & 0.089 \\
$n=12$ & 0.7421 & 15.932 & 1.321 & 2.000 & 0.195 & 0.173 \\
Pure delay & 0.7424 & 15.867 & 1.318 & 2.000 & 0.218 & 0.192 \\
\bottomrule
\end{tabular}
\caption{Joint fits with free ambient and history. Ambient offsets are in degrees Celsius; history fractions $\gamma$ multiply the initial command. Gain units are degrees Celsius per command percentage point. Every displayed step offset reaches the imposed $+2$ degree bound. These are sensitivity fits under declared bounds, not independently measured ambient temperatures or histories.}
\label{tab:nuisance-parameters}
\end{table}

% End expanded input


The nuisance fit retains the first sensor observation as its thermal initial
condition while changing the assumed ambient and delay-state history. Its
history fraction describes a declared state family; it does not reconstruct an
unrecorded heater command. The separate constant-input derivation in
Appendix~\ref{sec:nuisance-assumptions} explains why a fitted history or ambient
can also introduce exact parameter ambiguity. A favorable joint residual score
does not establish estimability in every experiment.

\begin{figure}[p]
 \centering
 \includegraphics[width=\linewidth]{nuisance_sensitivity.pdf}
 \caption{Ambient and initial-history assumptions change joint calibration
 error and the selected delay. Cold, fixed is the original linear-hold fit;
 History frees one delay-state history fraction per record with ambient fixed;
 Ambient frees one ambient offset per record with cold delay states; Both frees
 both. Offset bounds are $\pm2\,{}^\circ\mathrm C$ and history fractions lie in
 $[0,1]$. The two thermal experiments share their model parameters. Lines join
 categorical assumptions; markers are fitted representatives without confidence
 intervals. The one-stage delay uses the selected short-delay pole label.}
 \label{fig:nuisance-sensitivity}
\end{figure}

\subsection{Command reconstruction changes the inferred delay}
Changing the command hold shifts the fitted delay while producing smaller
changes in calibration RMSE (Fig.~\ref{fig:reconstruction-sensitivity}). A
backward-Euler fit can even reduce the residual score relative to exact
propagation. That does not make its discretization a more accurate description
of the continuous thermal system: fitted parameters can absorb numerical and
command-reconstruction differences. The recorded time/command pairs do not
identify the command waveform between samples. Consequently, the narrowest
conditional delay intervals should not be interpreted as resolving this
unmeasured waveform choice. All compared holds use the same observations and
the declared cold prehistory; pure delay has no backward-Euler counterpart in
this study.

\begin{figure}[p]
 \centering
 \includegraphics[width=\linewidth]{reconstruction_sensitivity.pdf}
 \caption{Different command reconstructions and propagators change the
 inferred delay at similar residual error. Each point is a separate joint cold
 fit. Linear, left and right denote the assumed hold between recorded command
 values. Exact uses the analytic LTI propagator for that hold; Right Euler uses
 the inherited backward-Euler/right-end convention at the recorded observation
 times. The pure-delay curve ends at the exact right-hold case. Lines join
 categorical assumptions and imply neither a convergence path nor uncertainty
 intervals.}
 \label{fig:reconstruction-sensitivity}
\end{figure}

\subsection{Prediction precision remains conditional on the mean model}
Prediction bands provide a coordinate-invariant target even when the original
physical coordinates have an exact null direction. Its zero-state equivalence
family preserves the measured thermal response, so prediction variance can be
computed on the identifiable quotient without assigning finite uncertainty to
each physical parameter. By contrast, constant-input calibration with
equilibrated delay states leaves a globally free delay that affects a varying
future input. Those target bands are withheld. The saved finite delay-exploration
envelopes are sensitivity plots without a probability interpretation.

For the jointly fitted cold representatives, mean-band widths depend strongly
on the assumed covariance or residual resampling
(Fig.~\ref{fig:prediction-mean-widths}). Fitted IID variance gives narrower mean
bands than the displayed HAC, OU and block-residual alternatives. Wider bands
under another working error law do not establish uncertainty from between-record
mean discrepancy. Table~\ref{tab:prediction-widths} separates mean-band width
from descriptive containment by observation bands on V01. Adding fitted future
error variance does not resolve that record's transfer error. V01 is the first
file in the existing sorted code map, with no selection by residual size; the
complete saved summary retains every validation file, including those with
better containment.

% Expanded input: tables/prediction_widths
% Generated by report_followups.py; scalar sources in followup_plot_data.json.
\begin{table}[tbp]
\centering\small
\begin{tabular}{lrrr}
\toprule
Model, method & \shortstack{Sine width\\(${}^\circ$C)} & \shortstack{V01 width\\(${}^\circ$C)} & V01 obs. fraction \\
\midrule
$n=1$, IID fitted & 0.078 & 0.061 & 0.135 \\
$n=1$, OU fitted & 0.439 & 0.347 & 0.139 \\
$n=2$, IID fitted & 0.088 & 0.077 & 0.140 \\
$n=2$, OU fitted & 0.544 & 0.479 & 0.154 \\
$n=12$, IID fitted & 0.099 & 0.094 & 0.150 \\
$n=12$, OU fitted & 0.662 & 0.629 & 0.166 \\
Pure delay, IID fitted & 0.102 & 0.097 & 0.161 \\
Pure delay, OU fitted & 0.560 & 0.512 & 0.171 \\
\bottomrule
\end{tabular}
\caption{Median pointwise mean-band widths and descriptive V01 containment by nominal observation-prediction bands. The last column uses observation bands with an added fitted IID or OU marginal future-error variance; it is not containment by the narrower mean bands. V01 is defined in Table~\ref{tab:validation-files}. These fractions on one observed record are not repeated-experiment coverage estimates. HAC and block-residual methods supply no future-error law here and therefore no observation-band containment.}
\label{tab:prediction-widths}
\end{table}

% End expanded input


\begin{figure}[p]
 \centering
 \includegraphics[width=\linewidth]{prediction_mean_widths.pdf}
 \caption{Conditional thermal-mean precision depends on the working error
 model. Points give median pointwise nominal 95\% mean-band widths over the
 sine calibration record and held-out V01, excluding their conditioned first
 observations. IID fitted uses residual degrees of freedom on the identifiable
 quotient; HAC60 uses Bartlett bandwidth 60; OU fitted uses the observed
 mean/length Hessian with variance profiled; Block60 uses percentile predictions
 from saved joint circular-block draws. The OU means are separately fitted;
 other methods retain the baseline mean. Lines join methods as visual guides,
 without a common unconditional confidence or repeated-experiment coverage
 interpretation. These are mean bands, not future-observation bands.}
 \label{fig:prediction-mean-widths}
\end{figure}

% End expanded input



% Expanded input: sections/design_results
\subsection{Local input design improves declared information criteria}
All parameter-design cases completed after the shutdown-grid correction, and
the independent audit passed the information matrices, eigenpairs, objective
values, selected-start rule and resource constraints. Each selected local
solution improves its own declared reference criterion. This comparison is
conditional on the fitted model points, working covariance and initialization;
it does not select a physically correct model.

The ensemble designs show how the target changes the proposed experiment
(Table~\ref{tab:parameter-design} and Fig.~\ref{fig:design-utilities}). The
minimum criterion targets the least favorable relative scenario; the mean
criterion targets average log improvement. Historical information changes both the objective
and the reference: future-only ensembles use multisine because constant input
leaves delay unobservable in equilibrated scenarios, whereas adding the two
historical calibration records makes the constant-reference total information
nonsingular. Differences between the two panels therefore cannot be read as
absolute information loss from including historical data. The determinant,
minimum-eigenvalue, delay and prediction criteria also have different scales.

% Expanded input: tables/parameter_design
% Generated by plot_design_results.py from separately audited phase summaries.
\begin{table}[tbp]
\centering\small
\begin{tabular}{llrr}
\toprule
Ensemble, information & Criterion & $\Delta\log$ utility & Max. (${}^\circ$C) \\
\midrule
mean, future & D & 1.508 & 84.990 \\
mean, history & D & 1.390 & 84.990 \\
mean, future & E & 1.423 & 74.333 \\
mean, history & E & 0.754 & 80.792 \\
mean, future & Delay & 1.439 & 74.333 \\
mean, history & Delay & 0.758 & 80.792 \\
mean, future & Prediction & 0.861 & 84.990 \\
mean, history & Prediction & 0.348 & 84.990 \\
min, future & D & 1.294 & 84.990 \\
min, history & D & 1.299 & 84.151 \\
min, future & E & 0.560 & 84.990 \\
min, history & E & 0.621 & 75.349 \\
min, future & Delay & 0.620 & 73.829 \\
min, history & Delay & 0.631 & 75.463 \\
min, future & Prediction & 0.704 & 84.990 \\
min, history & Prediction & 0.300 & 84.990 \\
\bottomrule
\end{tabular}
\caption{Selected local ensemble designs over 48 model/startup/covariance scenarios. Utility is the mean or minimum scenario-relative log improvement. Future-only ensembles use multisine as reference; ensembles with historical information use constant input, whose total information is nonsingular after adding calibration information. Criterion scales differ. Temperature is the largest independently checked value across sixteen fitted mean models through the shutdown tail. These local results do not establish global optimality, probabilistic model robustness or hardware limits.}
\label{tab:parameter-design}
\end{table}

% End expanded input


Figure~\ref{fig:design-commands} displays two selected ensemble determinant
designs and their dense temperature checks. The upper temperature envelope
ranges over the declared sixteen fitted mean models, without probabilities
over models or parameter uncertainty. The selected inputs use the amplitude
and slew limits and approach the imposed temperature constraint. The finite
tail check includes the last command's ramp to zero. All selected designs pass
the independent dense acceptance check, but this does not establish a
continuous-time or hardware temperature bound.

\begin{figure}[p]
 \centering
 \includegraphics[width=\linewidth]{design_commands_temperature.pdf}
 \caption{Information design changes the command and its predicted thermal
 envelope. Constant and multisine are common candidates; ensemble mean D and
 min D are selected determinant designs using future information only. Commands
 share the declared horizon, time-average dose, amplitude and slew constraints.
 The lower panel is the pointwise maximum temperature among sixteen fitted
 mean/startup combinations, not a confidence band or one physical trajectory.
 The vertical line marks the end of the design horizon and start of shutdown;
 the horizontal line is the assumed temperature acceptance limit. Calculations
 use a $0.25$ s audit grid through the $300$ s tail.}
 \label{fig:design-commands}
\end{figure}

\begin{figure}[p]
 \centering
 \includegraphics[width=\linewidth]{design_relative_utilities.pdf}
 \caption{The design target and scenario aggregate change relative utility.
 Points are the independently audited selected local objective values. Mean
 and min aggregate scenario-relative log improvements over 48 fixed
 model/startup/covariance scenarios. The first panel uses multisine as reference;
 the second adds historical calibration information and uses constant input.
 Criteria have different targets and scales, so ordinate heights across
 criteria or reference changes do not rank absolute physical information.
 Lines join categorical criteria as guides. No global optimality certificate
 is implied.}
 \label{fig:design-utilities}
\end{figure}

\subsection{Refitting competitors changes model discrimination}
The separate fixed-mean discrimination designs compare all six pairs among
the representative formulations (Table~\ref{tab:discrimination-design}). Their
mean/min objectives concern log improvements relative to multisine for each
pair. The smallest raw pair distance is a different summary and need not rank
those objectives in the same order. These designs share the parameter-design
resource assumptions but optimize mean separation rather than parameter
information.

% Expanded input: tables/discrimination_design
% Generated by plot_design_results.py from separately audited phase summaries.
\begin{table}[tbp]
\centering\small
\begin{tabular}{llrr}
\toprule
Startup, covariance & Aggregate & $\Delta\log D^2$ & Smallest $D^2$ \\
\midrule
Cold, IID & mean & 0.852 & 52.329 \\
Cold, IID & min & 0.798 & 52.329 \\
Cold, OU10 & mean & 0.825 & 17.797 \\
Cold, OU10 & min & 0.652 & 18.054 \\
Cold, OU60 & mean & 0.949 & 59.994 \\
Cold, OU60 & min & 0.628 & 58.797 \\
Equilibrated, IID & mean & 1.184 & 18.955 \\
Equilibrated, IID & min & 1.246 & 19.525 \\
Equilibrated, OU10 & mean & 1.451 & 2.136 \\
Equilibrated, OU10 & min & 1.324 & 1.564 \\
Equilibrated, OU60 & mean & 2.596 & 6.202 \\
Equilibrated, OU60 & min & 2.368 & 6.202 \\
\bottomrule
\end{tabular}
\caption{Fixed-mean discrimination designs compare six pairs among the four representative formulations. The objective aggregates pairwise log improvements relative to multisine, whereas the last column is the smallest raw selected pair distance; these are different comparisons. Covariances are common known working laws with marginal standard deviation $0.25\,{}^\circ\mathrm C$. Alternative-parameter refitting is evaluated separately in Fig.~\ref{fig:discrimination-refits}.}
\label{tab:discrimination-design}
\end{table}

% End expanded input


Allowing the competing formulation to change parameters reduces its distance
from the fixed one-stage generating mean
(Fig.~\ref{fig:discrimination-refits}). Under equilibrated constant input, the
closest-found distance vanishes to numerical precision for every alternative
and working covariance. The matched gain and time constant give an analytic
zero-distance witness because the delay factor is absent. Positive separation
at the original fitted alternative parameters had therefore overstated the
information in that input about model order.

The selected varying inputs retain larger closest-found distances than
multisine in this bounded directional comparison, including after refitting.
That result concerns a fixed one-stage generating mean against three
alternatives, rather than every possible generating formulation or parameter
point. Each smallest converged distance is an upper bound on the global closest
alternative. The ideal equal-prior Gaussian error formula for two fixed means
cannot be reused as a power calculation after fitting those alternatives or
allowing the observed between-record discrepancy.

\begin{figure}[p]
 \centering
 \includegraphics[width=\linewidth]{discrimination_refits.pdf}
 \caption{Parameter refitting can remove apparent model-order separation.
 The fixed one-stage generating mean is compared with the three indicated
 alternative formulations under the first declared covariance, IID Gaussian
 error with marginal standard deviation $0.25\,{}^\circ\mathrm C$. Rows use
 cold and input-equilibrated delay states. Left panels retain each alternative's
 original fitted parameters; right panels use its smallest converged
 24-start distance. Const. and Multi. are common candidates; Mean and Min are
 the selected fixed-mean discrimination inputs. Symmetric-log axes are linear
 below one, retaining vanishing distances without a positive plotting floor.
 Lines join categorical waveforms. The complete saved comparison also includes
 both correlated working covariances. Distances after refitting are upper
 bounds on global minima, without a real-data testing-power interpretation.}
 \label{fig:discrimination-refits}
\end{figure}

% End expanded input

\FloatBarrier

% End expanded input


% Expanded input: sections/discussion
\section{Discussion}
The transfer-function analysis places a limit on what a better input can achieve.
Under the single measured output and fixed initial conditions, excitation cannot
separate the original model's equivalent physical parameter combinations.
Estimating the three transfer coefficients or measuring an additional independent
quantity is a different scientific task from narrowing a four-parameter local
covariance. The existing data include a second sensor, but its usefulness for
breaking this equivalence requires a model that connects that sensor to the
states; the second trace alone does not establish that connection.

The reduced conductance--capacity parameters require an effective-model
interpretation. The gain fixes $\Ua=c/K$ only conditional on the adopted heater
conversion. The reduced heat capacity $\Cp=c\tau/K$ is tied to the selected
thermal pole; it is not automatically the heater heat capacity $\CpH$ of the
original model. The original model's equivalence class permits $\CpH$ to vary
while both measured transfer coefficients and temperature predictions are fixed.

Initialization is part of the model specification. A delay that appears estimable
under cold delay states may become unobservable under equilibrated delay states
and constant input. Future laboratory records should include the heater history
before the fitting window and an independent ambient measurement. These records
would reduce dependence on assumptions made at the first sample.

The conditional delay and physical-family-boundary intervals vary with residual
block length (Figs.~\ref{fig:delay-intervals} and~\ref{fig:family-boundaries}).
The circular bootstrap resamples centered residuals within each calibration
record and retains dependence within each block. The Gaussian comparator uses
a fixed hypothetical noise scale, whereas the block methods use empirical
centered residuals. Their contrast can therefore reflect both error scale and
serial dependence. Centering also
removes each calibration record's mean residual. This construction therefore
does not estimate the between-record offsets seen in validation, even when its
percentile intervals are wider than those from independent Gaussian simulation.
Selecting the widest of these conditional intervals would not supply a
probability model for the omitted discrepancy.

Fitting an OU correlation length
changes the weighting of thermal trajectories and their local uncertainty, but
the remaining negative innovation correlation challenges that combined
mean/covariance model. The fitted length cannot yet be interpreted as a measured
physical disturbance time scale. Neither a Gaussian likelihood improvement nor
a HAC variance calculation establishes that the observed error process satisfies
the assumptions needed for confidence coverage.

Transfer to the held-out files is a stronger concern than the small differences
among the representative cold model orders. The pooled error and its file-mean
decomposition (Table~\ref{tab:validation-summary}) show why narrower calibration
intervals alone would be insufficient evidence for using a formulation on a new
record. Subtracting each held-out mean error would use information unavailable
before that record was observed. Independent ambient, command-history and
heater-power measurements could help distinguish candidate explanations;
the supplied filenames and temperature offsets do not identify one.

The controlled coverage experiment isolates one consequence of a wrong error
law: IID likelihood intervals can fail even when the thermal mean is exactly
correct. HAC improves the comparison without reaching the nominal reference
at the longer tested correlation length. The finite-grid profile diagnostic
also prevents treating every interpolated threshold set as an exact numerical
confidence set. Neither result supplies coverage for the real-data bootstrap,
whose centering removes the between-record mean errors seen in validation.

The ambient/history and command-hold studies change a different part of the
inference. Ambient offsets lower calibration error but reach imposed bounds;
command reconstruction shifts the selected delay with less visible change in
RMSE. These sensitivities can exceed the apparent precision of a chosen
conditional error model. Recording independent ambient and heater history would
therefore answer a question that collecting more samples under the same
unverified initialization cannot settle.

The prediction calculation preserves the distinction between a redundant
physical representation and an unresolved future response. Quotient prediction
variance remains finite along the original physical equivalence family, whereas
an unobserved delay from a constant-input calibration can change a later varying
response. A pseudoinverse alone cannot decide which of those situations applies.
The prospective calculations specify what a proposed experiment is meant to
learn. Parameter-information and model-separation criteria choose different
inputs; initialization and the working covariance remain part of each design.
Refitting the alternatives removes constant-input separation in the equilibrated
case, while the selected varying inputs retain separation in the bounded
one-stage-generating comparison. None of those designs can remove the original
physical equivalence family under the same modeled output.

The next laboratory study should record an independent ambient trace, the
pre-experiment command history and heater-power calibration, then test a
selected varying input against held-out predictions declared beforehand. A
second sensor could address the physical ambiguity only after its observation
equation is specified and checked. Those measurements would test explanations
for the transfer errors that the current conditional intervals and fixed-point
input designs leave unresolved.

% End expanded input


% Expanded input: sections/provenance
\section{Data, code, and research status}
The repository source is \url{https://github.com/dowlinglab/pyomo-doe}.
The study is under \path{skunk_works/delay_study} on branch
\path{reformulation-exploration}. Phase manifests identify immutable source
commits, data hashes, configuration, environment, seeds, and numerical outcomes.
Release receipts identify the source revision, compiled report checksum, numerical phase manifests and artifact checks. The supplied raw data and inherited source were preserved.

An OpenAI Codex agent performed the notebook exports, implemented the numerical
analysis, ran the computations, searched and inspected literature, and drafted
this document. This disclosure describes scientific and software work as well
as writing assistance. The document has not received human scientific review.
No laboratory hardware experiment was performed for this study, and no target
journal or submission policy has been selected.

% End expanded input

\bibliographystyle{plainnat}
\bibliography{bibliography}
\appendix

% Expanded input: sections/structural_derivations
\section{Structural derivations}
\label{sec:structural}

\subsection{Delay-chain moments and pure-delay limit}
Let $\lambda_n=n/\theta$ be the delay-state rate. The normalized chain kernel for
$t\geq0$ is
\[
 h_n(t)=\frac{\lambda_n^n t^{n-1}e^{-\lambda_n t}}{(n-1)!}.
\]
Its Laplace transform is $(1+s/\lambda_n)^{-n}$. Differentiating that transform at zero
gives mean $n/\lambda_n=\theta$ and variance $n/\lambda_n^2=\theta^2/n$. Consequently, the
normalized complete thermal response has impulse-time mean $\tau+\theta$ and
variance $\tau^2+\theta^2/n$. These are properties of the model kernel;
they are not estimates of a measured residence-time distribution.
For fixed $s$ in the domain of convergence,
\[
 \log(1+\theta s/n)^{-n}=-n\log(1+\theta s/n)\longrightarrow-\theta s.
\]
This calculation supplies the pure-delay limit of Eq.~\eqref{eq:transfer}.

\subsection{An explicit physical-parameter equivalence class}
Define the two positive pole time constants $t_{\mathrm f}<t_{\mathrm s}$ by
$a=t_{\mathrm f}+t_{\mathrm s}$ and $b=t_{\mathrm f}t_{\mathrm s}$. For fixed
positive transfer coefficients $(K,a,b)$, choose a family coordinate
$r=\CpS/\Ub$ between these pole time constants. The physical parameters
\begin{equation}
 \Ua=\frac{c}{K},\qquad
 \CpH=\frac{b\Ua}{r},\qquad
 \CpS=\Ua\left(a-r-\frac{b}{r}\right),\qquad
 \Ub=\frac{\CpS}{r}
 \label{eq:equivalent}
\end{equation}
produce exactly the same transfer function for every
$r\in(t_{\mathrm f},t_{\mathrm s})$. Positivity follows from
$a-r-b/r=-(r-t_{\mathrm f})(r-t_{\mathrm s})/r>0$. Substitution into
Eq.~\eqref{eq:original-transfer} recovers $K$, $a$, and $b$ exactly.
The conductance $\Ua$ is identifiable through $K$; the other three parameters
vary along this equivalence class. No input waveform can distinguish members
of this class through this output under the stated initial conditions.

Positivity constrains this continuum to finite ranges in some coordinates.
For example, its heat capacities obey
\[
 \Ua t_{\mathrm f}<\CpH<\Ua t_{\mathrm s},\qquad
 0<\CpS\leq\Ua\left(a-2\sqrt b\right).
\]
The second upper bound occurs at $r=\sqrt b$.
The coupling conductance satisfies
\[
 0<\Ub\leq\Ua\left(\frac{a^2}{4b}-1\right),
\]
with its maximum at $r=2b/a$.
Thus bounded physical ranges or bounded exploration grids do not establish
structural identifiability. Structural uniqueness must be checked separately
from the numerical width of a confidence set.

The one-stage reduced model has $a=\tau+\theta$ and $b=\tau\theta$.
For a fixed transfer function, its two pole labels therefore assign reduced
capacities $\Cp=\Ua\tau$ equal to $\Ua t_{\mathrm f}$ and
$\Ua t_{\mathrm s}$. These values are precisely the infimum and supremum of
the original heater capacity, rather than interior members of its strictly
positive family. A capacity obtained from either reduced pole label cannot
be identified with a unique original $\CpH$. The boundary functions themselves
are identifiable from $(K,a,b)$ and their sampling uncertainty can be propagated
without selecting a family coordinate $r$.

The strictly positive original parameters imply distinct poles. To see this,
write $h=\CpH/\Ua$, $r=\CpS/\Ub$, and $q=\CpS/\Ua$. Then $a=h+r+q$
and $b=hr$, giving $a>h+r\geq2\sqrt b$. Coincident poles are attained only
as a limiting case outside this strictly positive four-parameter family.

\subsection{Constant input with equilibrated delay states}
If $u(t)=u_0$ and $z_i(0)=u_0$, every delay-state derivative in
Eq.~\eqref{eq:chain} vanishes. The predicted temperature then depends on
$K$ and $\tau$ alone, and its derivative with respect to $\theta$ is zero at
every observation time. A constant prehistory gives the same conclusion for
Eq.~\eqref{eq:puredelay}. This experiment cannot estimate the delay even when
the thermal transient is well measured. Cold delay states under the same
post-start input retain a transient that depends on $\theta$.

\subsection{Coordinate changes and information}
For positive parameters, write $\bm\eta=(\log K,\log\tau,\log\theta)^\top$
and $\bm\xi=(\log\Ua,\log\Cp,\log\theta)^\top$. The Jacobian of the map
from $\bm\eta$ to $\bm\xi$ is the constant matrix
\[
 M=\begin{pmatrix}-1&0&0\\-1&1&0\\0&0&1\end{pmatrix},\qquad M^{-1}=M.
\]
For a nonsingular covariance $C_{\eta}$, the transformed local covariance is
$C_{\xi}=M C_{\eta}M^\top$. Sensitivity rank and prediction-space uncertainty
are invariant under this invertible transformation. Numerical eigenvalues,
eigenvectors, and condition numbers generally change with the coordinate metric.
The matrix satisfies $|\det M|=1$, so log-coordinate information determinants
at the same model point are equal in this particular transformation.
This equality applies to future-input information and to any historical
information transformed by the same congruence. It therefore preserves the
D-design objective and its input gradient. The delay target is unchanged
because the third coordinate is unchanged, and a prediction loss is preserved
when its target matrix transforms with the sensitivities. The minimum
information eigenvalue does not have this invariance.

\subsection{Repeated poles and a singular local information matrix}
At a repeated thermal and delay pole, $\tau=\theta/n$, the transfer function
is $(1+\tau s)^{-(n+1)}K$. The derivatives in unscaled coordinates satisfy
\[
 \frac{\partial\log G_n(s)}{\partial\tau}
 =-\frac{s}{1+\tau s},\qquad
 \frac{\partial\log G_n(s)}{\partial\theta}
 =-\frac{s}{1+\theta s/n}.
\]
These two derivatives coincide at the repeated pole. The corresponding
log-coordinate sensitivity columns are proportional with ratio $n$, producing
a singular local information matrix under zero-state excitation.
For $n>1$, the rational transfer function still uniquely specifies the distinct
single and repeated poles away from their coincidence; at coincidence there is
one common pole value. Local derivative rank loss at this exceptional point
therefore does not by itself establish a continuum of exactly equivalent
parameters. Nonregular local inference and structural nonuniqueness are separate
questions. The exact matrix-exponential simulator avoids a partial-fraction
formula that would divide by the pole separation.

The exceptional point also changes the order of local objective curvature.
Let $t_{\mathrm c}>0$ be the common pole time and perturb
$\tau=t_{\mathrm c}+\delta$ and $\theta=nt_{\mathrm c}-\delta$, holding the
total mean time fixed. For small $\delta$,
\[
 \log\frac{G_n(s;\delta)}{G_n(s;0)}
 =\frac{n+1}{2n}\,\delta^2
   \left(\frac{s}{1+t_{\mathrm c}s}\right)^2+O(\delta^3).
\]
The linear term cancels. The temperature difference is therefore second order
in this direction and its squared prediction distance is fourth order, subject
to a nonzero response coefficient for the chosen input. For $n=1$, exchanging
$\delta$ and $-\delta$ exchanges pole labels exactly. For $n>1$, the repeated
factor's multiplicity generally distinguishes those two perturbed models.

\subsection{Moment coordinates}
Define the total mean response time $m=\tau+\theta$ and the delay fraction
$f=\theta/m$. The inverse map is $\tau=m(1-f)$ and $\theta=mf$ with
$m>0$ and $0<f<1$. If $\ell=\log[f/(1-f)]$, the Jacobian from
$(\log m,\ell)$ to $(\log\tau,\log\theta)$ is
\[
 \begin{pmatrix}1&-f\\1&1-f\end{pmatrix}.
\]
This determinant-one map provides an alternative optimization metric without
changing model predictions. Its rectangular bounds generally differ from the
original rectangular bounds in $\tau$ and $\theta$.

The impulse-time variance $v=\tau^2+\theta^2/n$ provides another summary of the
transfer shape. Given $m$ and $v$, the possible delay values solve
\[
 \left(1+\frac1n\right)\theta^2-2m\theta+(m^2-v)=0.
\]
Two positive solutions can therefore share the first two moments. For $n=1$
they exchange the two poles and give the same full transfer function. For
$n>1$, higher moments can distinguish the branches. The minimum feasible
variance, $v=m^2/(n+1)$, occurs at the repeated-pole condition above. These
identities motivate comparing coordinate conditioning and response moments,
while retaining the full dynamic model for estimation.

\subsection{A measurement that could break the original equivalence}
Suppose an additional sensor measured the model's heater state $T_{\mathrm H}$
directly. Let $Y_{\mathrm H}(s)$ and $Y_{\mathrm S}(s)$ denote the Laplace
transforms of the heater and sensor temperature deviations from ambient.
For zero initial deviations, the second balance in Eq.~\eqref{eq:original} gives
\[
 \frac{Y_{\mathrm H}(s)}{Y_{\mathrm S}(s)}
 =1+\frac{\CpS}{\Ub}s=1+rs,
\]
for nonzero transfer responses. Identifying this additional time constant $r$ would select one
member of Eq.~\eqref{eq:equivalent} and hence recover all four physical
parameters. The supplied
second TC-Lab temperature trace has not been established as a direct
measurement of this model's heater state.

\subsection{Dependent-residual sandwich calculation}
For a full-rank sensitivity matrix, let $B=(J^\top J)^{-1}$ and define the
score contribution $\bm q_t=J_t^\top e_t$ at residual index $t$ within an
experiment. With lag bandwidth $L$, the matrix used between the two factors
of $B$ is
\[
 S=\sum_t\bm q_t\bm q_t^\top+
 \sum_{j=1}^{L}\left(1-\frac{j}{L+1}\right)
 \sum_{t=j+1}^{N}
 \left(\bm q_t\bm q_{t-j}^\top+\bm q_{t-j}\bm q_t^\top\right).
\]
For joint calibration we formed each lagged sum within its experiment, then
added the experiment contributions; no lag joined separate experiments. The
reported covariance is $[N/(N-k)]BSB$, where $k$ is the parameter count and
$N$ is the total conditional residual count. The scalar degrees-of-freedom
adjustment is an implementation choice. Bartlett weighting follows the
positive-semidefinite construction of \citet{newey-1987-simple}; it does not
establish that the TC-Lab score process satisfies the paper's consistency
assumptions. Lags here count retained observations rather than elapsed seconds.

% End expanded input


% Expanded input: sections/nuisance_assumptions
\section{Calibration, initial-state and residual assumptions}
\label{sec:nuisance-assumptions}

\subsection{Heater calibration introduces another exact ambiguity}
The conversion $c=\alpha P$ is fixed in the baseline study. If it is unknown,
the transformation
\[
 (c,\Ua,\Cp)\longmapsto(\lambda c,\lambda\Ua,\lambda\Cp),
 \qquad \lambda>0,
\]
preserves $K=c/\Ua$, $\tau=\Cp/\Ua$, and every reduced-model prediction under
the same initial conditions. In the original model, scaling $c,\Ua,\Ub,\CpH$
and $\CpS$ together has the same effect. Input design cannot identify this
absolute thermal scale from temperature and percentage command alone.

Let $C_{\mathrm{io}}$ be a covariance matrix for
$(\log K,\log\tau,\log\theta)$ and suppose an independent heater calibration
has log-variance $v_c$. With the transformation matrix $M$ defined in
Appendix~\ref{sec:structural}, the corresponding physical log-covariance is
\[
 C_{\mathrm{phys}}=M C_{\mathrm{io}}M^\top+v_c\bm v\bm v^\top,
 \qquad \bm v=(1,1,0)^\top.
\]
This covariance identity is exact in logarithmic coordinates under the stated
independence assumption; Gaussianity is not needed for the covariance itself.
It adds correlated uncertainty to conductance and capacity while leaving delay
unchanged. A hypothetical lognormal calibration with coefficient of variation
$r_c$ has $v_c=\log(1+r_c^2)$. Such a sensitivity calculation must be labelled
as assumed calibration uncertainty: these data do not supply $r_c$.

\subsection{Ambient and history are nuisance quantities}
For either prescribed delay-state initialization, hold the observed initial
sensor temperature $T_0$ fixed and vary ambient temperature. Linearity gives
\[
 \frac{\partial T(t)}{\partial\Tamb}=1-e^{-t/\tau}.
\]
An ambient error therefore persists rather than averaging away with additional
samples in one record. Under constant input $u_0$ and equilibrated delay states,
\[
 T(t)=T_0 e^{-t/\tau}+(\Tamb+Ku_0)(1-e^{-t/\tau}).
\]
Only the combination $\Tamb+Ku_0$ is identified from this response if ambient
is free, and delay is absent. Independent ambient information and a varying
input address different ambiguities.

Another restricted history family sets every chain state initially to
$\gamma u_0$, where $0\leq\gamma\leq1$. For fixed parameters, its response is
exactly $(1-\gamma)T_{\mathrm{cold}}+\gamma T_{\mathrm{equilibrated}}$.
The same interpolation describes a constant pre-window input between zero and
$u_0$ in the pure-delay model. This family does not represent every possible
heater history. Its derivative with respect to $\gamma$ is the difference of
the two endpoint trajectories. If these nuisance sensitivities resemble the
delay sensitivity, fixing history can make delay appear more precise than a
fit that allows history to vary.

More generally, let $A=W^{1/2}J_p$ and $B=W^{1/2}J_a$ be the mean sensitivities
for parameters of interest and nuisance parameters, respectively, under a fixed
positive-definite observation weight $W$. Profiling unrestricted local nuisance
directions leaves information $A^\top(I-BB^+)A$, where $B^+$ is the
Moore--Penrose pseudoinverse. This projection cannot increase information.
Bounds or independent nuisance measurements alter the problem and must be stated.

\subsection{Free history can preserve the exchanged pole labels}
For the one-stage model, write $y=T-\Tamb$, $A_0=T_0-\Tamb$, and
$z(0)=\gamma u_0$, where $u_0=u(0)$. Taking the one-sided Laplace transform
of the two state equations gives
\[
 (1+\tau s)(1+\theta s)Y(s)
 =K U(s)+K\theta\gamma u_0+\tau A_0(1+\theta s).
\]
Exchange $\tau$ and $\theta$ while keeping $K$, ambient and the observed
initial temperature fixed. The denominator and the coefficient of $s$ in the
initial-state term stay the same. Equality of the remaining constant terms
requires, for $u_0\ne0$,
\begin{equation}
 \gamma'=\frac{\theta\gamma+A_0(\tau-\theta)/(K u_0)}{\tau}.
 \label{eq:history-pole-exchange}
\end{equation}
This transformation preserves the entire trajectory for any subsequent input.
It is admissible only if $\gamma'\in[0,1]$ and both parameter vectors satisfy
their fit bounds. For a joint fit, each record needs its own admissible
transformed history fraction. When ambient has an offset, $A_0$ uses that
record's actual fitted ambient.

At $A_0=0$, equation~\eqref{eq:history-pole-exchange} reduces to
$\gamma'=\gamma\theta/\tau$. A prescribed equilibrated history can separate
pole labels that become indistinguishable when history is free. If $u_0=0$,
history fractions have no effect in this restricted family; the exchange
requires $A_0(\tau-\theta)=0$ instead. A local full-rank sensitivity matrix
cannot exclude an admissible disconnected pole-label solution.

\subsection{Constant input admits a gain--ambient--history continuum}
For a constant command $u_0$, define the delay-state deficit
\[
d_n(t)=e^{-nt/\theta}\sum_{k=0}^{n-1}\frac{(nt/\theta)^k}{k!}
\]
for an $n$-stage chain, and $d_\infty(t)=\mathbf 1_{\{t<\theta\}}$
for pure delay. With the uniform initial history fraction $\gamma$, the
thermal forcing is $K u_0[1-(1-\gamma)d_n(t)]$. Thus
\[
 T(t)=T_0 e^{-t/\tau}+C(1-e^{-t/\tau})
 -\frac{B u_0}{\tau}\int_0^t e^{-(t-v)/\tau}d_n(v)\,dv,
 \qquad C=\Tamb+K u_0,\quad B=K(1-\gamma).
\]
The same expression applies to pure delay with $d_\infty$. Holding $\tau$,
$\theta$, $B$ and $C$ fixed gives the exact family
\begin{equation}
 K'>0,\qquad
 \gamma'=1-\frac{K(1-\gamma)}{K'},\qquad
 \Tamb'=\Tamb+u_0(K-K').
 \label{eq:constant-nuisance-gauge}
\end{equation}
Every admissible member predicts the same measured-temperature trajectory.
The gain box, ambient-offset bounds and $0\leq\gamma'\leq1$ restrict the
family. At an interior point they still permit a continuum. With
$\gamma=1$, the delay is also absent from the response. Varying the input
generally breaks equation~\eqref{eq:constant-nuisance-gauge}, since the
time-varying gain forcing cannot be absorbed into a constant ambient offset.
Independent state-propagation checks verify this family for one-, two- and
twelve-stage chains and pure delay, and the pole exchange for a varying input.

\subsection{Serial dependence alone does not prove least-squares bias}
For fixed input and known initial conditions, suppose observations have the
correct mean $\bm m(\bm p_0)$ and additive error $\bm\epsilon$ with
$E\bm\epsilon=0$ and covariance $\Sigma$. For any candidate parameter vector
$\bm p$,
\[
 E\|\bm m(\bm p)-\bm m(\bm p_0)-\bm\epsilon\|^2
 =\|\bm m(\bm p)-\bm m(\bm p_0)\|^2+\operatorname{tr}\Sigma.
\]
Correlated errors change sampling uncertainty and efficiency without moving the
minimum of this expected output-error objective. This identity does not imply
an unbiased finite-sample nonlinear estimator. It also does not apply unchanged
when noisy observations enter regressors, the initial observation creates a
nonzero conditional error mean, or the thermal mean is misspecified. The
observed TC-Lab residual offsets and held-out errors require those possibilities
to remain open.

\subsection{An OU filter can anticorrelate an IID component}
Consider a fixed OU correlation factor $0<\phi<1$ at regular sampling and an
IID component $v_i$ with variance $s_v^2$. Away from the first retained
innovation, the unnormalized filtered component is $w_i=v_i-\phi v_{i-1}$.
Independence gives
\[
 \operatorname{Var}(w_i)=(1+\phi^2)s_v^2,\qquad
 \operatorname{Cov}(w_i,w_{i-1})=-\phi s_v^2,\qquad
 \rho_w(1)=-\frac{\phi}{1+\phi^2}.
\]
Dividing by the constant OU innovation scale does not change this correlation.
It approaches $-1/2$ as $\phi\to1$. Thus an IID component inside a residual
process can produce negative filtered correlation even when slower components
are present. This calculation concerns one component and a fixed regular-sample
filter. It neither establishes that the TC-Lab has this decomposition nor
assigns the observed correlations to sensor noise rather than mean error.

\section{Chain convergence at the causal onset}
\label{sec:onset}
For a cold constant step of magnitude $u_0$ and zero initial thermal deviation,
let $D_n$ have an Erlang
distribution with mean $\theta$ and variance $\theta^2/n$. The temperature
rise relative to ambient is
\[
 y_n(t)=K u_0 E\!\left[
  \left(1-e^{-(t-D_n)/\tau}\right)\mathbf 1_{\{D_n\leq t\}}\right].
\]
At $t=\theta$, the pure-delay response is zero while every finite chain has a
positive response. The central limit theorem and a first-order expansion of
the exponential give
\[
 \sqrt n\,y_n(\theta)\longrightarrow
 \frac{K u_0\theta}{\tau\sqrt{2\pi}}.
\]
To justify the expectation limit, the standardized Erlang variables have unit
second moments, so their positive parts are uniformly integrable. The
exponential remainder is bounded by a constant times
$E(\theta-D_n)^2/\tau^2=O(n^{-1})$.
This is an onset error of order $n^{-1/2}$, even when smooth parts of the
response converge faster. A small average trajectory error can therefore hide
a larger difference near the delay's causal onset. The saved onset calculation
checks the coefficient using an independently tested Erlang-CDF expression.

% End expanded input


% Expanded input: sections/frequency_information
\section{Frequency sensitivity and local excitation requirements}
\label{sec:frequency-information}
The following identities are derived for the transfer functions in this study.
They explain an ideal periodic steady-state experiment. They do not replace
the finite-horizon, initialized simulations or their working covariance models.
For positive frequency $\omega$, let $G=G_n(i\omega)$ and define the complex
log-transfer sensitivity vector $\bm h=G^{-1}\partial G/\partial\bm\eta$,
where $\bm\eta=(\log K,\log\tau,\log\theta)$. Direct differentiation gives
\begin{equation}
 \bm h_n(\omega)=\left(
 1,\;\frac{-i\omega\tau}{1+i\omega\tau},\;
 \frac{-i\omega\theta}{1+i\omega\theta/n}
 \right)^\top.
 \label{eq:frequency-sensitivities}
\end{equation}
For the pure-delay model, the third component is $-i\omega\theta$.
At zero frequency, only the gain sensitivity remains. At low frequency,
\[
 \log G_n(i\omega)=\log K-i\omega(\tau+\theta)
 -\tfrac12\omega^2(\tau^2+\theta^2/n)+O(\omega^3).
\]
Thus the leading phase change depends on the sum of the thermal and delay
lags. Their separation depends on higher-order frequency behavior; a small
low-frequency fit error can coexist with poor precision for each lag.
The $\theta^2/n$ term is absent in the pure-delay expansion.

To separate this local frequency geometry from startup information, consider
a single periodic input $u(t)=u_0+A\cos\omega t$ after transients have decayed.
Let $\bm d=\partial G(i\omega)/\partial\bm\eta$ and suppose the $N$ observed
phases have zero sine/cosine sums, zero sine--cosine cross sum, and equal
sine/cosine sums of squares $N/2$.
With gain direction $\bm e_1=(1,0,0)^\top$ and independent errors of known
variance $\sigma^2$, the exact phase-sum information is
\begin{equation}
 F=\frac{N}{\sigma^2}\left[
 (u_0K)^2\bm e_1\bm e_1^\top+
 \frac{A^2}{2}\operatorname{Re}(\bm d\bm d^{*})\right],
 \label{eq:periodic-information}
\end{equation}
where $*$ denotes conjugate transpose. A single tone without its constant (DC) offset
contributes a matrix of rank at most two for the three real parameters.
With nonzero offset and amplitude, the DC term supplies a separate gain direction:
its constant sensitivity is orthogonal to the oscillatory components over these
phases. The determinant that tests independence of the lag
components is
\begin{equation}
 \begin{vmatrix}
 \operatorname{Re}h_\tau&\operatorname{Re}h_\theta\\
 \operatorname{Im}h_\tau&\operatorname{Im}h_\theta
 \end{vmatrix}
 =\frac{\omega^3\tau\theta(\tau-\theta/n)}
 {(1+\omega^2\tau^2)(1+\omega^2\theta^2/n^2)}.
 \label{eq:frequency-lag-area}
\end{equation}
It is nonzero away from $\tau=\theta/n$. In this ideal experiment, a DC
component and one nonzero tone therefore give local full rank at a regular
parameter point. For pure delay the determinant is
$\omega^3\tau^2\theta/(1+\omega^2\tau^2)>0$.
At the finite-chain repeated pole, $h_\theta=nh_\tau$ at every frequency;
adding frequencies cannot remove that first-order rank deficiency.
These are local differential statements. They do not remove the cold one-stage
global pole-label ambiguity, calibrate a confidence interval, or establish an
optimal input under amplitude, sampling and temperature constraints.

The factor $G$ in $\bm d=G\bm h$ also matters: information depends on output
sensitivity, rather than phase sensitivity alone. The time-domain designs
evaluate that output sensitivity with retained observations, startup assumptions
and a specified covariance. Constant-command startup records can carry transient
information that is deliberately absent from the periodic steady-state calculation.
Analytic derivatives, the determinant and the phase-sum information were checked
against parameter perturbations, numerical determinants and explicit real
sensitivity matrices. These checks contain no additional TC-Lab fits.

% End expanded input


% Expanded input: sections/literature_comparison
\section{Comparison of theoretical assumptions}
\label{sec:literature-comparison}
The review covers inspected passages from seven cited works. Two library
carriers supply the profile paper and a continuous-time identification chapter;
five further primary PDFs were inspected outside the ingested library. This
counts works and carriers, without implying that every page was read. The
source and search records retain exact queries, dates, versions, page locations
and access gaps. Table~\ref{tab:literature-comparison} separates the inspected
theory from its use in this study.

\begingroup
\small
\begin{longtable}{>{\raggedright\arraybackslash}p{.17\linewidth}>{\raggedright\arraybackslash}p{.25\linewidth}>{\raggedright\arraybackslash}p{.25\linewidth}>{\raggedright\arraybackslash}p{.19\linewidth}}
 \caption{Questions, source assumptions and the scope of transfer to TC-Lab.
 Page numbers are physical PDF pages in the inspected source. The last column
 gives this study's interpretation rather than a result established by the
 cited work.}\label{tab:literature-comparison}\\
 \toprule
 Question and source & Inspected theory & Conditions or distinction & TC-Lab use and limit \\
 \midrule
 \endfirsthead
 \toprule
 Question and source & Inspected theory & Conditions or distinction & TC-Lab use and limit \\
 \midrule
 \endhead
 \bottomrule
 \endfoot
 What produces a residual?\newline
 \citet{garnier-2008-direct}, pp.~27--29
 & The input-driven response and residual process are separate descriptions.
 Residuals can combine measurement error, disturbances and model inadequacy.
 & The chapter provides an identification framework, without selecting a
 correct residual mechanism for these experiments.
 & Treat covariance fits as working assumptions; examine transfer to other
 records separately. \\[.6em]

 Can parameters be separated?\newline
 \citet{raue-2009-structural}, pp.~1--5
 & Profile likelihood fixes a parameter and reoptimizes nuisance parameters;
 its shape can differ from local quadratic curvature.
 & A likelihood and a numerical exploration are required. A sampled crossing
 alone cannot certify every nuisance minimum or mode.
 & Combine profiles with exact transfer equivalence and sensitivity rank;
 retain disconnected components. \\[.6em]

 What should information target?\newline
 \citet{hjalmarsson-2009-system}, pp.~8--9
 & Application loss and confidence regions can concern useful parameter
 combinations rather than uniform precision in every coordinate.
 & The inspected prediction-error information formula assumes the true system
 is in the model set and the predictor is differentiable.
 & Target delay or thermal prediction while checking mean-model adequacy and
 the chosen parameter metric. \\[.6em]

 How is a bounded design optimized?\newline
 \citet{manchester-2010-input}, pp.~1--5
 & Finite-horizon amplitude-constrained information design is nonconvex; the
 source develops a convex relaxation.
 & The inspected source is the eight-page author preprint. Its relaxation
 guarantees concern its own construction and assumptions.
 & Our multistart nonlinear solutions have no inherited approximation or
 global-optimality guarantee. \\[.6em]

 What does a chain represent?\newline
 \citet{hurtado-2019-generalizations}, pp.~11--12
 & Erlang dwell times admit a first-order differential-equation chain, with
 equivalence to an integral description.
 & The source treats mean-field transitions. Our constant-rate input kernel
 uses the corresponding Erlang construction.
 & Derive kernel moments and limits; a fitted stage count alone does not
 establish a thermal transport mechanism. \\[.6em]

 How can serial dependence affect covariance?\newline
 \citet{newey-1987-simple}, pp.~1--3
 & Bartlett-weighted long-run covariance estimation is positive semidefinite;
 consistency has explicit score, moment and dependence conditions.
 & Asymptotic bandwidth conditions are distinct from choosing several fixed
 bandwidths for a finite dataset.
 & Compare HAC sensitivity without treating it as proof of coverage or a
 correction for mean-model bias. \\[.6em]

 What dependence does resampling retain?\newline
 \citet{politis-1994-stationary}, pp.~1--3
 & Periodic extension supports circular resampling. Remark~2 describes the
 fixed-length relative of the geometric stationary-bootstrap scheme.
 & The schemes differ. Equal marginal weights at record edges do not establish
 nonlinear interval coverage with model discrepancy.
 & Vary fixed block lengths and retain within-record centering; keep the
 bootstrap uncertainty conditional. \\
\end{longtable}
\endgroup

The sources address different parts of the same inference problem. A
continuous-time response model defines the quantities that can be identified;
an information objective then assigns precision to local directions or
predictions. A profile explores beyond that local quadratic approximation.
The original TC-Lab model's exact equivalent family remains indistinguishable
under every input in the stated experiment, so selecting a larger determinant
cannot resolve its physical ambiguity.

Covariance estimation and resampling also answer different questions. HAC
uses residual score products to estimate estimator variability; a circular
bootstrap generates new records from resampled residual blocks. Both require
assumptions about the errors around the fitted mean. A systematic error on a
different experiment is evidence about mean transfer and experimental
conditions, rather than evidence that either covariance procedure has measured
that discrepancy.

The next measurement questions follow from these limits. An independent
heater-power calibration would address the absolute conductance--capacity
scale. Ambient and command prehistory records would constrain startup nuisance
assumptions. A measurement equation connecting the second sensor to a modeled
state is needed before using it to select an original physical family member.
The adaptive-design journal paper and the
unavailable bootstrap source listed in the search record remain reading leads,
without theorem claims assigned to unread versions.

% End expanded input


% Expanded input: sections/source_audit
\section{Literature coverage and claim audit}
The review used the configured lab literature skill and project-oriented searches
of the local library, followed by primary-source web discovery. It is a focused
review, without an exhaustive database search or a formal systematic-review
protocol. Search passages that only list another work in their references were
excluded as support for that work's claims.

The source record \path{literature/source_review.json} retains paper identities,
source events where available, physical-page locators, and claim limitations.
The record \path{literature/acquired_sources.json} identifies downloaded PDFs
by their hashes. Full paper text and PDFs remain outside this repository.

For \citet{raue-2009-structural}, the library paper identifier is
\texttt{6925c25f}\ldots\texttt{ae8160}; the inspected profile construction is
on physical PDF page 4, printed page 1926, Eq.~(10). We checked that equation
against the rendered page. The journal source by
\citet{hjalmarsson-2009-system} was checked at physical pages 8--9, printed
pages 282--283, for confidence
regions, identifiable combinations, and the assumptions behind prediction-error
information. Those two pages were also inspected as rendered images because
the extraction engines disagree on equation symbols. The journal source by
\citet{hurtado-2019-generalizations} was checked at physical pages 11--12,
printed pages 1841--1842,
for the chain equations and their integral-equation equivalence, with rendered
page inspection of the mathematical passages.

The journal PDF of \citet{newey-1987-simple} was an image-only scan. We read physical pages
1--3, printed pages 703--705, with optical character recognition and checked
the rendered pages. Equation (5) on physical page 2 gives the Bartlett weights;
Theorem 2 on page 3 states assumptions for consistency. The scanned equation
symbols were checked visually rather than copied from the recognition output.

The chapter by \citet{garnier-2008-direct} was inspected at physical pages 27--29 of the
stored book, printed pages 1--3. Its byline includes Peter C. Young in addition
to the two book editors. The residual decomposition is on physical page 28,
printed page 2. Chapter and editor roles are retained separately in the reference.
Physical-to-printed page offsets were checked at each passage rather than
extrapolated to the whole supplied PDF. Contextual chapter readings remain
separate leads in the source record.

The journal scan of \citet{politis-1994-stationary} was read with optical character recognition and
all three cited pages were checked as rendered images. Physical pages 1--3 are
printed pages 1303--1305. The circle convention is specified on physical page 1;
the fixed-length relative is discussed in Remark 2 on physical page 3. The
geometric-block stationarity and consistency results are not ascribed to our
fixed-length nonlinear residual bootstrap.

For \citet{manchester-2010-input}, we inspected the eight-page author preprint. We cite its
identifier, \path{arXiv:1009.5614}, and retain the distinct six-page conference
version as an unread publication record.
Its input-information formulation and optimization assumptions were inspected
in the body. The downloaded adaptive-input-design preprint changes title and
author order relative to its published journal record. Access to the journal
body remains unresolved, and it is not cited for a theorem. These version and
access limitations remain in the source records.

The exact expanded manuscript, citation inventory and claim assessments are
bound in the private evidence packet. Scoped public receipts record the draft
hash, inspected source identities, available fine-evidence bindings and explicit
gaps for externally inspected PDFs that have no ingested carrier. Two cited
works have ingested carriers; five primary PDFs were inspected externally.
Source availability, quote fidelity, claim support and human scientific review
remain separate statuses. Mechanical checks and agent assessments do not
constitute human acceptance of the manuscript.

% End expanded input

\clearpage

% Expanded input: sections/teaching_doe
\section{Teaching with reduced models and native Pyomo.DoE}
\label{sec:teaching-doe}
This supplement adds the teaching comparison and native Pyomo.DoE demonstration
completed after the overnight study. Its four designs use a separate,
single-model protocol from the ensemble designs in the main results. The saved
fits, uncertainty calculations and literature claims in the preceding sections
are unchanged.

\subsection{One stage is sufficient for the core teaching exercise}
Use $n=1$ to teach parameter estimation, local information and input design.
It has two dynamic states and three unknown transfer parameters, $K$, $\tau$
and $\theta$. It retains the original four-parameter model's response under the
studied observation equation and startup conditions while removing a redundant
parameter direction. Add $n=2$ as an extension on delay shape, model comparison
and numerical discretization. It has three states and still three unknown
parameters. The available fits do not require higher orders for this exercise.

The independently optimized joint $n=1$ and original fits differ in predicted
temperature by less than $1.5\times10^{-7}\,{}^\circ\mathrm C$ on calibration
and held-out records. Parameter reduction therefore preserves these predictions.
Table~\ref{tab:teaching-prediction} gives the cold-start comparison on the joint
calibration and eleven duplicate-excluded held-out files, without refitting
to those files. The $n=2$ joint calibration RMSE is about $12.5\%$ higher and
its pooled held-out RMSE about $0.06\%$ higher. On the step-only calibration,
$n=2$ lowers RMSE from $0.1573$ to $0.1529\,{}^\circ\mathrm C$, about $2.8\%$.
Initialization can reverse the order ranking. These results do not establish a
general predictive advantage for an additional delay stage.

% Expanded input: tables/teaching_prediction
% Generated by teaching_followup/render_report_tables.py
\begin{table}[htbp]
\centering
\small
\caption{Cold-start joint fits and pooled held-out prediction error. RMSE is in degrees Celsius.}
\label{tab:teaching-prediction}
\begin{tabular}{lrr}
\toprule
Model & Calibration RMSE & Held-out RMSE \\
\midrule
Original four-parameter & 0.4115 & 3.3588 \\
$n=1$ & 0.4115 & 3.3588 \\
$n=2$ & 0.4630 & 3.3608 \\
\bottomrule
\end{tabular}
\end{table}

% End expanded input


The $n=1$ pole labels can be exchanged globally. A convention such as
$\tau\geq\theta$ chooses the short-delay representative; a full-rank local
information matrix does not detect the other mode. The $n=2$ repeated delay pole
distinguishes it from a distinct thermal pole, while coincident poles remain
nonregular. Neither reduced formulation identifies a unique member of the
original physical parameter family.

\subsection{Residual correlation changes uncertainty more than model order}
Table~\ref{tab:teaching-blocks} compares nominal $95\%$ percentile intervals
from $400$ refits per method using the same centered empirical residual pool.
Changing independent resampling to $60$-sample circular blocks widens the delay
interval by factors of $6.66$ for $n=1$ and $6.23$ for $n=2$. Sixty samples
are roughly one minute in these records. This isolates retained dependence
from the change of marginal error scale that occurs when comparing empirical
resampling with fixed-scale Gaussian simulation.

% Expanded input: tables/teaching_blocks
% Generated by teaching_followup/render_report_tables.py
\begin{table}[htbp]
\centering
\small
\caption{Delay percentile intervals (s) from the same centered residual pool; 400 refits per method. Width ratio compares 60-sample blocks with independent resampling.}
\label{tab:teaching-blocks}
\begin{tabular}{lrrr}
\toprule
Model & Independent & 60-sample blocks & Width ratio \\
\midrule
$n=1$ & [16.41, 16.70] & [15.62, 17.56] & 6.66 \\
$n=2$ & [15.22, 15.56] & [14.31, 16.44] & 6.23 \\
\bottomrule
\end{tabular}
\end{table}

% End expanded input


These are conditional intervals: centering removes the record-specific mean
error, block length is an assumption, and the $n=1$ interval excludes its global
swapped-pole mode. The controlled OU60 simulation provides a separate coverage
check. With the thermal mean correct by construction, nominal $95\%$
IID fitted-variance delay intervals cover only $66/300$ datasets for $n=1$
and $56/300$ for $n=2$. HAC60 raises these counts to $235/300$ and $249/300$.
The real-data one-component OU fits retain innovation lag-one correlations near
$-0.5$. Neither the wider bootstrap intervals nor the OU likelihood establishes
an adequate real-data discrepancy model.

\subsection{A three-parameter Pyomo.DoE implementation}
The demonstration uses Pyomo $6.10.0$ and a Parmest \texttt{Experiment} subclass
with the experiment-input, experiment-output, measurement-error and
unknown-parameter suffixes. The implementation follows the
\href{https://pyomo.readthedocs.io/en/stable/explanation/analysis/doe/guide.html}{official Pyomo.DoE quick-start guide}
(accessed October 21, 2026). Each reduced model yields a full-rank $3\times3$
local information matrix. Under the original four-parameter single-output
formulation, the exact four-dimensional determinant vanishes because the rank
is three. Finite-difference roundoff cannot restore the missing information.

Let $p=(K,\tau,\theta)$, $\eta=\log p$ and
$J_{ij}=\partial\mu(t_i)/\partial\eta_j$. Relative-parameter scaling in
Pyomo.DoE supplies this local log-parameter sensitivity. The design maximizes
the log determinant of
\begin{equation}
 F_\eta=J^\mathsf{T}\Sigma^{-1}J.
\end{equation}
There is no prior information matrix. The nominal parameters are the saved joint
cold fits. Each of the four model/covariance cases uses two optimization starts,
constant $50\%$ and alternating $30/70\%$ command. The selected local result
has the larger converged objective; no global optimality claim follows.

The experiment lasts $900$ s and observes temperature every $10$ s, excluding
the initial value, for $90$ observations. Seven command knots are spaced
$150$ s apart, with linear interpolation. Commands lie between $0$ and $100\%$,
adjacent knots differ by at most $60$ percentage points, and integrated command
is at most $45000$ percent-seconds. The ambient and initial temperature are
$22.2\,{}^\circ\mathrm C$, with cold delay states. Nominal collocation
temperatures are constrained below $84.99\,{}^\circ\mathrm C$, retaining a
$0.01\,{}^\circ\mathrm C$ guard below the independent $85\,{}^\circ\mathrm C$
acceptance limit.

The measurement-error suffix contains standard deviations. Its built-in
weighting assumes independent errors. For the correlated comparator, the
stationary OU law has known marginal standard deviation
$\sigma=0.25\,{}^\circ\mathrm C$ and known correlation length $60$ s.
With $\phi_i=\exp[-(t_i-t_{i-1})/60]$, transformed mean outputs are
\begin{equation}
 \widetilde\mu_1=\mu_1,\qquad
 \widetilde\mu_i=\frac{\mu_i-\phi_i\mu_{i-1}}{\sqrt{1-\phi_i^2}},
 \quad i\geq2.
\end{equation}
Applying independent weighting with standard deviation $\sigma$ to these
outputs gives the full-covariance information matrix above. Covariance
parameters are held fixed throughout.

\subsection{The selected inputs heat, cool and heat again}
Table~\ref{tab:teaching-doe} compares each selected input with constant $50\%$
under the same model and working covariance. The D-information multiplier is
$[\det(F_{\mathrm{selected}})/\det(F_{\mathrm{constant}})]^{1/3}$.
It ranges from $3.67$ to $4.12$. The local delay SD falls from
$0.728$ to $0.272$ s for $n=1$ under independent errors and from
$1.627$ to $0.678$ s under known OU60 errors. The corresponding $n=2$ values
are $0.630$ to $0.245$ s and $1.225$ to $0.536$ s.

% Expanded input: tables/teaching_doe
% Generated by teaching_followup/render_report_tables.py
\begin{table}[htbp]
\centering
\small
\caption{Conditional local delay standard deviations (s) and D-information multipliers for the selected and constant-input designs. Each row uses its own working covariance and a hypothetical marginal error SD of $0.25\,{}^\circ\mathrm C$.}
\label{tab:teaching-doe}
\begin{tabular}{llrrr}
\toprule
Model & Error law & Constant SD & Selected SD & D multiplier \\
\midrule
$n=1$ & IID & 0.728 & 0.272 & 4.12 \\
$n=2$ & IID & 0.630 & 0.245 & 4.09 \\
$n=1$ & Known OU60 & 1.627 & 0.678 & 3.80 \\
$n=2$ & Known OU60 & 1.225 & 0.536 & 3.67 \\
\bottomrule
\end{tabular}
\end{table}

% End expanded input


These SDs concern hypothetical future measurements at the fixed fitted
parameters and error law. They are separate from the empirical residual
intervals above. The information gains do not measure a realized improvement
in prediction accuracy.

Figure~\ref{fig:teaching-experiments} shows the selected commands and their
nominal temperature responses. All four inputs start at $100\%$ for $2.5$ min,
ramp down through about $86\%$ at $5$ min and $26$--$34\%$ at $7.5$ min,
and reach zero at $10$ min. They remain mostly off until $12.5$ min, then ramp
to $60$--$66\%$ at $15$ min. Each uses the same integrated command as the
constant $50\%$ reference, concentrating heat early to create heating and
cooling transients. Table~\ref{tab:teaching-commands} supplies the knot values.

The predicted peak is about $84.93\,{}^\circ\mathrm C$ at $5.4$--$5.5$ min.
Changing the working covariance has a modest effect on the $n=1$ schedule.
For $n=2$, the two selected commands differ by less than
$5\times10^{-7}$ percentage points. Thus nearly identical input schedules can
coexist with substantially different uncertainty calculations. The correlated
assumption gives selected-input delay SDs about $2.2$--$2.5$ times their
independent-error counterparts in this comparison.

\begin{figure}[p]
 \centering
 \includegraphics[width=\linewidth]{../teaching_followup/figures/optimized_experiments.png}
 \caption{Native Pyomo.DoE local D-optimal teaching experiments for three
 transfer parameters. Columns compare one and two delay stages. Upper panels
 show affine heater commands; lower panels show separately simulated nominal
 temperature responses on a $0.25$ s grid. Blue curves assume independent
 errors; orange curves assume known stationary OU correlation of $60$ s.
 Gray dotted curves use constant $50\%$ command with the same total dose.
 The red line is the imposed $85\,{}^\circ\mathrm C$ acceptance limit.
 Curves are conditional cold-start means; they are not uncertainty bands.}
 \label{fig:teaching-experiments}
\end{figure}

% Expanded input: tables/teaching_commands
% Generated by teaching_followup/render_report_tables.py
\begin{table}[htbp]
\centering
\small
\caption{Selected heater command knots (percent), rounded to two decimals. Commands interpolate linearly; exact values are saved in the command CSV.}
\label{tab:teaching-commands}
\begin{tabular}{rrrrr}
\toprule
Time (min) & $n=1$ IID & $n=1$ OU60 & $n=2$ IID & $n=2$ OU60 \\
\midrule
0.0 & 100.00 & 100.00 & 100.00 & 100.00 \\
2.5 & 100.00 & 100.00 & 100.00 & 100.00 \\
5.0 & 85.71 & 85.60 & 86.06 & 86.06 \\
7.5 & 25.71 & 34.40 & 33.94 & 33.94 \\
10.0 & 0.00 & 0.00 & 0.00 & 0.00 \\
12.5 & 5.71 & 0.00 & 0.00 & 0.00 \\
15.0 & 65.71 & 60.00 & 60.00 & 60.00 \\
\bottomrule
\end{tabular}
\end{table}

% End expanded input

\FloatBarrier

\subsection{Numerical checks and classroom scope}
An independent auditor checked all sixteen saved information matrices and local
SD evaluations: the baseline, two starts and a refined-mesh evaluation in each
case. It uses complex-step sensitivities from the study's exact simulator and
explicit full covariance inversion, independently of the Pyomo whitening
implementation. The largest relative FIM error is below $3.3\times10^{-7}$;
the largest relative local-SD error is below $2.6\times10^{-5}$.

Three-point Radau collocation uses $90$ finite elements for $n=1$ and $180$
for $n=2$, with the selected inputs also evaluated at twice those counts.
The initial $n=2$ mesh failed the $0.001\,{}^\circ\mathrm C$ mean tolerance
and was refined. An unnecessary lower bound on the collocation temperature
polynomial caused false infeasibility near startup; removing that bound
preserved the upper temperature constraint. Failed trials and the correction
to sensitivity evaluation at active design constraints are retained in
\path{teaching_followup/trial_failures.json}.

All four selected commands satisfy dose, amplitude and knot-change checks.
Their independently simulated temperatures remain below $85\,{}^\circ\mathrm C$
on the $0.25$ s grid over the proposed $900$ s window. This check covers the
nominal fitted points. Parameter uncertainty, continuous-time bounds, shutdown
after $900$ s and hardware operation require separate assessment.

For teaching, fit $n=1$, expose the original physical redundancy, compare
independent and correlated uncertainty, and optimize a varying input.
Use $n=2$ to ask whether an added delay stage changes prediction, global
identifiability or numerical requirements. The scripts, saved results,
independent audit and plotting receipt are under
\path{skunk_works/delay_study/teaching_followup}.

% End expanded input

\end{document}
